\providecommand{\toplevelprefix}{../..}  %
\documentclass[../../book-main.tex]{subfiles}

\begin{document}

\chapter{Pursuing Low-Dimensional Distributions via Denoising}
\label{ch:denoising}

\begin{quote}
    \hfill ``\textit{Information is the resolution of uncertainty}.''

    \hfill ---Claude Shannon, 1948
\end{quote}
\vspace{5mm}

In \Cref{ch:classic-models}, we showed how to learn simple classes of
distributions whose supports are assumed to be either a single
low-dimensional subspace, a mixture of such subspaces, or low-rank
Gaussians. For simplicity, the different (hidden) linear or Gaussian
modes are taken to be orthogonal or independent,\footnote{Or can be
reduced to such idealistic cases under benign conditions.} as
illustrated in \Cref{fig:subspaces}. For these special distributions,
simple and effective learning algorithms with correctness and efficiency
guarantees can be derived, and the geometric and statistical
interpretations of the associated algorithmic operations are clear.

In practice, both linearity and independence are idealistic
assumptions that distributions of real-world high-dimensional data
\textit{rarely} satisfy. The only assumption we can make is that the
intrinsic dimension of the distribution is very low compared with the
dimension of the ambient space in which the data are embedded. Hence,
in this chapter we show how to learn a more general class of
low-dimensional distributions in a high-dimensional space that is
\textit{not} necessarily (piecewise) linear.

Real-world distributions typically contain multiple components or
modes, for example, corresponding to different object classes in
images. These modes may not be statistically independent and can even
have different intrinsic dimensions. Therefore, we generally
assume that the data are distributed on a mixture of low-dimensional
nonlinear submanifolds embedded in a high-dimensional space.
\Cref{fig:mixture-manifolds} illustrates such a distribution. Another
constraint is that in real-world situations we typically have access
only to a finite number of samples from the data distribution,
rather than prior knowledge of the exact distribution itself.

To learn such a distribution under these conditions, we must address
three fundamental questions:
\begin{itemize}
    \item What is a general approach to learning and representing a
        low-dimensional distribution in a high-dimensional space?
    \item How do we measure the complexity of the distribution so
        that we can exploit its low dimensionality effectively for learning?
    \item How do we make the learning process computationally
        tractable and scalable, especially when the ambient
        dimension is high and the sample size is large?
\end{itemize}
As we will see, the fundamental idea of \textit{compression}, or
\textit{dimension reduction}, which we have already shown to be effective in the
linear/independent case, remains a guiding principle for developing
computational models and methods for general nonlinear
low-dimensional distributions.

Due to its theoretical and practical importance, we examine in greater
depth how this compression-based framework materializes when the
distribution of interest can be well modeled or approximated by a
mixture of low-dimensional subspaces or low-rank Gaussians. As later
chapters show, these seemingly restrictive model classes serve as basic
building blocks for algorithms that perform compression on general
distributions.

\begin{figure}
    \centering
    \includegraphics[width=0.5\linewidth]{\toplevelprefix/chapters/chapter3/figs/mixed-manifolds.png}
    \caption{Data \(\vX\) are distributed on a mixture of
        low-dimensional submanifolds \(\cup_{j} \cM_{j}\) in a very
        high-dimensional ambient space, say \(\R^{D}\). This is a much
        more general case than the special cases studied in
        \Cref{ch:classic-models}, where the submanifolds were assumed to
        be piecewise linear and thus admitted parametric forms and
    analytical solutions.}
    \label{fig:mixture-manifolds}
\end{figure}

\section{Entropy Minimization and Compression}

\subsection{Entropy and Coding Rate}
In \Cref{ch:intro} we mentioned that the goal of learning is to
find the simplest way to generate a given set of data. Conceptually,
Kolmogorov complexity was intended to provide such a measure of
complexity, but it is not computable and is not associated with any
implementable scheme that can actually reproduce the data. Hence, we
need an alternative, computable, and realizable measure of
complexity. This leads us to the notion of \textit{entropy},
introduced by Shannon in 1948~\cite{Shannon-1948}.

To illustrate the constructive nature of entropy, let us start with
the simplest case. Suppose we have a discrete random variable
that takes \(N\) distinct values, or \textit{tokens}, \(\{\vx_{1},
\dots, \vx_{N}\}\) with equal probability \(1/N\). We can
encode each token \(\vx_{i}\) using the \(\log_{2} N\)-bit binary
representation of \(i\). This coding scheme can be generalized to
encoding arbitrary discrete distributions~\cite{Cover-Thomas}: given
a distribution \(p\) such that \(\sum_{i=1}^{N} p(\vx_{i}) = 1\), one
can assign each token \(\vx_{i}\) with probability \(p(\vx_{i})\)
a binary code of length \(\log_{2} [1/p(\vx_{i})] = - \log_{2}
p(\vx_{i})\) bits. Hence the average number of bits, or \textit{the
coding rate}, needed to encode any sample from the distribution
\(p(\cdot)\) is given by the expression\footnote{By the convention
    of information theory~\cite{Cover-Thomas}, the \(\log\) here has
base \(2\). Hence, entropy is measured in (binary) bits.}:
\begin{equation}
    H(p) \doteq \Ex_{\vx \sim p}[-\log p(\vx)] = - \sum_{i=1}^{N}
    p(\vx_{i}) \log p(\vx_{i}).
    \label{eqn:entropy-discrete}
\end{equation}
This is known as the \textit{entropy} of the (discrete) distribution
\(p(\cdot)\). Note that this entropy is always nonnegative and it is
zero if and only if \(p(\vx_{i}) = 1\) for some \(\vx_{i}\) with \(i
\in [N]\).\footnote{Here we use the fact
\(\lim_{p \to 0} p \log p = 0\).} For a random variable
\(\vx\), we will often write \(H(\vx)\) to mean the entropy of the
(marginal) distribution of \(\vx\).

\subsection{Differential Entropy}

When the random variable \(\vx \in \R^{D}\) is continuous and has a
probability density \(p\), one may view the limit of the above
sum \eqref{eqn:entropy-discrete} as related to an integral:
\begin{equation}
    h(p) \doteq \Ex_{\vx \sim p}[-\log p(\vx)] = - \int_{\R^{D}}
    p(\vxi) \log p(\vxi) \odif{\vxi}.
    \label{eqn:entropy-differential}
\end{equation}
More precisely, given a continuous probability distribution \(p\) on
\(\R^{D}\), we may discretize \(\R^{D}\) into a union of disjoint
cubes \(C_{i}^{\epsilon}\) of side length \(\epsilon^{1/D}\), each (say)
centered at \(\vx_{i}\), and define the discretized distribution
\(p^{\epsilon}\) by \(p^{\epsilon}(\vx_{i}) = \Pr_{\vx \sim p}[\vx
\in C_{i}^{\epsilon}]\). Then one can show that \(H(p^{\epsilon}) +
\log \epsilon \approx h(p)\). By taking \(\epsilon\) to be small, we
observe that the differential entropy \(h(p)\) can be negative.
Interested readers may refer to \cite{Cover-Thomas} for a more
detailed explanation. Similar to the discrete entropy, we will often
write \(h(\vx)\) to mean the entropy of the (marginal) distribution of \(\vx\).

\begin{example}[Entropy of Gaussian distributions]
    Through direct calculation, one can show that the entropy of a
    Gaussian random variable \(x \sim \dNorm(\mu, \sigma^{2})\) is
    \begin{equation}
        h(x) = \frac{1}{2}\log(2\pi\sigma^{2}) + \frac{1}{2}.
        \label{eqn:entropy-Gaussian}
    \end{equation}
    The entropy of a multivariate Gaussian random variable \(\vx \sim
    \dNorm(\vmu, \vSigma)\) in \(\R^{D}\) is
    \begin{equation}
        h(\vx) = \frac{D}{2}(1 + \log(2\pi)) +
        \frac{1}{2}\logdet(\vSigma).
        \label{eqn:entropy-Gaussian-multi}
    \end{equation}
\end{example}

Similar to the entropy of a discrete distribution, we would like the
differential entropy to be associated with the coding rate of some
realizable coding scheme. For example, as mentioned above, if we
discretize the domain of the distribution with a grid of size
\(\epsilon > 0\) with centers \(\vx_{i}\), one possible coding scheme
is to map each \(\vx\) to its nearest center \(\vx_{i}\), and
the entropy of the resulting discrete distribution approximates the
differential entropy \cite{Cover-Thomas}.

There are some caveats associated with the definition
of differential entropy. For a distribution in a high-dimensional
space, when its support becomes degenerate (low-dimensional), its
differential entropy diverges to \(-\infty\). This fact is proved in
\Cref{thm:max_entropy}, but even in the simple explicit case of
Gaussian distributions \eqref{eqn:entropy-Gaussian-multi}, when the
covariance \(\vSigma\) is singular, we can see that
\(\logdet(\vSigma) = -\infty\) so we have \(h(\vx) = -\infty\) (note
    that by \Cref{thm:max_entropy}, Gaussian distributions have the
largest entropy among all distributions with the same mean and covariance). In
such a situation, it is not obvious how to properly quantize or
encode such a distribution. Nevertheless, degenerate (Gaussian)
distributions are precisely the simplest possible, and arguably the
most important, instances of low-dimensional distributions in a
high-dimensional space. In the next chapter, we will discuss a complete
resolution to this seeming difficulty with degeneracy.

\subsection{Minimizing Coding Rate}\label{sub:min_entropy}
The learning problem entails recovering a (potentially continuous)
distribution \(p(\vx)\) from samples \(\{\vx_{1}, \dots, \vx_{N}\}\)
drawn from it. For ease of exposition, we write \(\vX = \mat{\vx_{1},
\dots, \vx_{N}} \in \R^{D \times N}\). Because the distributions of
interest are (nearly) low-dimensional, we expect their (differential)
entropy to be very small. However, unlike the situations studied in
the previous chapter, we do not know the family of (analytical)
low-dimensional models to which \(p(\vx)\) belongs. Thus, checking
whether the entropy is small seems to be the only guideline we can
rely on to identify and model the distribution.

Now, given the samples alone without knowing what \(p(\vx)\) is, in
theory they could be interpreted as samples from any generic
distribution. In particular, they could be interpreted as any of the following:
\begin{enumerate}
    \item as samples from the empirical distribution \(p^{\vX}\)
        itself, which assigns probability \(1/N\) to each of the
        \(N\) samples \(\vx_{i}, i \in [N]\);
    \item as samples from a standard normal distribution \(\vx^{n}
        \sim p^{n} \doteq \dNorm(\vzero, \sigma^{2} \vI)\) with
        variance \(\sigma^{2}\) large enough (say larger than the
        squared sample norms);
    \item as samples from a normal distribution \(\vx^{e} \sim p^{e}
        \doteq \dNorm(\vzero, \hat{\vSigma})\) with covariance
        \(\hat{\vSigma} = \frac{1}{N} \vX \vX^{\top}\) being the
        empirical covariance of the samples;
    \item as samples from a distribution \(\hat{\vx} \sim
        \hat{q}(\vx)\) that closely approximates the ground-truth
        distribution \(p\).
\end{enumerate}
The question is: which interpretation is better, and in what sense? Suppose
that you believe these data \(\vX\) are drawn from a particular
distribution \(q(\vx)\), which may be one of the above. Then we could
encode the data points with the optimal codebook for the distribution
\(q(\vx)\). Then, as the number of samples \(N\) becomes large, the
required average coding length (or coding rate) is given by the
so-called \textit{cross-(differential)-entropy}:
\begin{equation}
    \CE(p \mmid q) \doteq \Ex_{\vx \sim p}[-\log q(\vx)].
\end{equation}
If we have identified
the correct distribution \(p(\vx)\), the coding rate is given by the
entropy \(h(p) = \Ex_{\vx \sim p}[-\log p(\vx)]\). It turns out
that the above coding rate \(\CE(p \mmid q)\) is always greater than
or equal to the entropy \(h(\vx)\) unless \(p = q\).

Hence, given a set of sampled data \(\vX\), to determine which
distribution best characterizes the data among
\(p^{n}\), \(p^{e}\), and \(\hat{q}\), we could encode the data \(\vX\)
using the codebooks associated with each of these distributions,
and compare their coding rates to see which yields the lowest rate:
\begin{equation}
    \argmin_{q} \CE(p \mmid q)
\end{equation}

However, since the true distribution \(p\) is unknown,
we must look for an alternative metric to evaluate.
For these specific choices of models—\(p^{n}\), \(p^{e}\), and \(\hat{q}\)—
their coding rates are exactly or nearly identical to their own differential entropies.
Since their entropies decrease progressively:
\begin{equation}
    h(p^{n}) > h(p^{e}) > h(\hat{q}).
\end{equation}
it naturally follows that the coding rates for \(\vX\) decrease in the same order:
\begin{equation}
    \CE(p \mmid p^{n}) > \CE(p \mmid p^{e}) > \CE(p \mmid \hat{q}).
\end{equation}

This observation gives a general guideline for pursuing a
distribution \(p(\vx)\) that has low-dimensional structure. It
suggests two possible approaches:
\begin{enumerate}[label={A\arabic*.}, leftmargin=*]
    \item Starting with a general distribution (say a normal
        distribution) with high entropy, gradually transform the
        distribution toward the (empirical) distribution of the data
        by reducing entropy;
    \item Among a large family of (parametric or non-parametric)
        distributions with explicit coding schemes that encode the
        given data, progressively search for better coding schemes
        that give lower coding rates.
\end{enumerate}
Conceptually, both approaches A1 and A2 try to do the same thing.
For A1, we need to ensure such a path of transformation exists and is
computable. For A2, it is necessary that the chosen family is rich
enough and can closely approximate (or contain) the ground-truth
distribution. For either approach, we need to ensure that solutions
with lower entropy or better coding rates can be efficiently computed
and converge to the desired distribution quickly.\footnote{Say the
distribution of real-world data such as images and texts.} We will
explore both approaches in the remaining sections of this chapter.

\section{Compression via Denoising}\label{sub:compression_denoising}

In this section we describe a \textit{natural} and
\textit{computationally tractable} way to learn a distribution
\(p_{\vx}\) by learning a parametric encoding that minimizes the
entropy or coding rate of the representation, then using this encoding
to transform high-entropy samples from a standard Gaussian into
low-entropy samples from the target distribution, as illustrated in
\Cref{fig:diffusion-chapter3}.  This methodology combines approaches
A1 and A2 above to learn and sample from the distribution.

\begin{figure}[t]
    \centering
    \includegraphics[width=\linewidth]{\toplevelprefix/chapters/chapter3/figs/diffusion_pipeline.png}
    \caption{Illustration of an iterative denoising process that,
        starting from an isotropic Gaussian distribution, converges to an
    arbitrary data distribution.}
    \label{fig:diffusion-chapter3}
\end{figure}

\subsection{Diffusion and Denoising Processes}
\label{sub:intro_diffusion_denoising}

We first want a procedure that reduces the entropy of a very noisy
random variable, producing a lower-entropy sample from the data distribution.
Here we describe one potential approach, perhaps the most natural.
First, we devise a way to \textit{gradually increase} the entropy of
samples already drawn from the data distribution.  Next, we seek an
\textit{approximate inverse} of this process.  In general, increasing
entropy has no inverse, because information from the original
distribution may be destroyed.  We therefore restrict ourselves to a
special case in which (1) the entropy-increasing operation is simple,
computable, and reversible, and (2) we have a (parametric) encoding of
the data distribution, as alluded to above.  As we will see, these two
conditions make our approach feasible.

\subsubsection{Diffusion Processes}

We will increase the entropy in arguably the simplest possible way,
i.e., \textit{adding isotropic Gaussian noise}. More precisely, given
the random variable \(\vx\), we can consider the \textit{stochastic
process} \((\vx_{t})_{t \in [0, T]}\) that adds gradual noise to it, i.e.,
\begin{equation}\label{eq:additive_gaussian_noise_model}
    \vx_{t} \doteq \vx + t\vg, \qquad \forall t \in [0, T],
\end{equation}
where \(T \in [0, \infty)\) is a time horizon and \(\vg \sim
\dNorm(\vzero, \vI)\) is drawn independently of \(\vx\). This process
is an example of a \textit{diffusion process}, so named because it
spreads (i.e., \textit{diffuses}) the probability mass over all
of \(\R^{D}\) as time goes on, increasing the entropy. This
intuition is confirmed graphically by \Cref{fig:ve_forward_density},
and rigorously via the following theorem.
\begin{theorem}[Simplified version of \Cref{thm:diffusion_entropy_increases}]
    Suppose that \((\vx_{t})_{t \in [0, T]}\) follows the model
    \eqref{eq:additive_gaussian_noise_model}. For any \(t \in (0,
    T]\), the random variable \(\vx_{t}\) has a probability density
    \(p_{t}\) and differential entropy
    \(h(\vx_{t}) > -\infty\). Moreover, under certain technical
    conditions on \(\vx\),
    \begin{equation}
        \odv*{h(\vx_{t})}{t} > 0, \qquad \forall t \in (0, T],
    \end{equation}
    i.e., the entropy of the distribution of \(\vx_{t}\) increases
    over time \(t\).
\end{theorem}
The proof is elementary, but it is rather long, so we postpone it to
\Cref{sub:diffusion_entropy_increases}. The main as-yet unstated
implication of this result is that \(h(\vx_{t}) > h(\vx_{s})\) for every
\(t > s \geq 0\). To see this, note that if \(s = 0\) and \(h(\vx) =
h(\vx_{0}) = -\infty\) then \(h(\vx_{t}) > -\infty\) for all \(t >
0\), so the conclusion is true; otherwise \(h(\vx_{t}) = h(\vx_{s}) +
\int_{s}^{t} \bs{\odv*{h(\vx_{u})}{u}} \odif{u} > h(\vx_{s})\) by the
fundamental theorem of calculus, so in both cases \(h(\vx_{t}) >
h(\vx_{s})\) for every \(t > s \geq 0\).

\begin{figure}
    \includegraphics[width=\textwidth]{\toplevelprefix/chapters/chapter3/figs/ve_forward_diffusion_density.png}
    \caption{\small\textbf{Diffusing a mixture of Gaussians \(\vx\).}
        From left to right, we observe the evolution of the density
        \(p_{t}\) of \(\vx_{t}\) as \(t\) grows from \(0\) to \(10\),
        along with some representative samples (red). The plane is
        colored by the probability density \(p_{t}\); high-density
        regions are colored darker blue. We observe that the probability
        mass becomes less concentrated as \(t\) increases, signaling that
    entropy increases.}
    \label{fig:ve_forward_density}
\end{figure}

\begin{example}[Adding noise to a mixture of
    Gaussians]\label{ex:gaussian_mixture}
    In this example, we consider adding noise to a mixture of \(K\)
    Gaussians. Choose \(K\) means
    \(\vmu_{1}, \dots, \vmu_{K} \in \R^{D}\), \(K\) covariance
    matrices \(\vSigma_{1}, \dots, \vSigma_{K} \in \PSD(D)\), and
    probabilities \(\pi_{1}, \dots, \pi_{K} \in [0, 1]\) such that
    \(\sum_{k = 1}^{K}\pi_{k} = 1\). Then \(\vx\) is sampled as
    follows.
    \begin{itemize}
        \item First, an index (or \textit{label}) \(y \in [K]\) is
            sampled s.t.~\(y = k\) with probability \(\pi_{k}\).
        \item Second, \(\vx\) is sampled from the Gaussian
            \(\dNorm(\vmu_{y}, \vSigma_{y})\); i.e., if \(y = k\) then
            \(\vx \sim \dNorm(\vmu_{k}, \vSigma_{k})\).
    \end{itemize}
    Thus \(\vx\) has the Gaussian mixture distribution
    \begin{equation}
        \vx \sim \sum_{k = 1}^{K}\pi_{k}\dNorm(\vmu_{k}, \vSigma_{k}).
    \end{equation}
    Gaussian mixture models are relevant because they can approximate
    many distributions: given any distribution for \(\vx\), there
    exists a Gaussian mixture that approximates it arbitrarily well.
    For this reason, Gaussian mixture models are a popular,
    analytically tractable surrogate for multi-modal and complex data
    distributions. In this book we use them extensively as a ``model
    problem'' to understand and develop almost every concept. For now,
    we represent the data \(\vx\) by a Gaussian mixture. After adding
    noise, the distribution of \(\vx_{t}\) is again a Gaussian
    mixture:
    \begin{equation}
        \vx_{t} = \vx + t\vg \sim \sum_{k =
        1}^{K}\pi_{k}\dNorm(\vmu_{k}, \vSigma_{k} + t^{2}\vI).
    \end{equation}
    To see by example how the entropy evolves, suppose \(K = 1\), so
    \(\vx \sim \dNorm(\vmu, \vSigma)\). By
    \eqref{eqn:entropy-Gaussian-multi}, the entropy of \(\vx_{t}\) is
    \begin{equation}
        h(\vx_{t}) = \frac{D}{2}(1 + \log(2\pi)) +
        \frac{1}{2}\logdet(\vSigma + t^{2}\vI).
    \end{equation}
    Differentiating with respect to time (see \Cref{app:optimization})
    gives
    \begin{equation}
        \odv*{h(\vx_{t})}{t} = t\tr((\vSigma + t^{2}\vI)^{-1}) =
        \sum_{i = 1}^{D}\frac{t}{\lambda_{i}(\vSigma) + t^{2}} > 0,
    \end{equation}
    where \(\lambda_{i}(\vSigma)\) are the (non-negative) eigenvalues of
    \(\vSigma\). Because this derivative is positive, the entropy
    \(h(\vx_{t})\) of the noised distribution increases over time.
\end{example}

\subsubsection{Denoising Processes}
The inverse operation to adding noise is known as \textit{denoising}.
It is a classical and well-studied topic in signal processing and
system theory, instantiated (for example) as the Wiener filter and
the Kalman filter (see \Cref{sec:analytical}). Several problems discussed in
\Cref{ch:classic-models}, such as PCA, ICA, and Dictionary Learning, are
specific instances of the denoising problem. For a fixed \(t\) and
the additive Gaussian noise model
\eqref{eq:additive_gaussian_noise_model}, the denoising problem can
be formulated as attempting to learn a function \(\bar{\vx}^{\star}(t,
\cdot)\) that forms the best possible approximation (in expectation)
of the true random variable \(\vx\), given both \(t\) and \(\vx_{t}\):
\begin{equation}\label{eq:denoising_loss}
    \bar{\vx}^{\star}(t, \cdot) \in \argmin_{\bar{\vx}(t,
    \cdot)} \Ex \norm{\vx - \bar{\vx}(t, \vx_{t})}_{2}^{2}.
\end{equation}
The solution to this problem, when optimizing \(\bar{\vx}(t, \cdot)\)
over all possible (square-integrable) functions, is the so-called
\textit{Bayes optimal denoiser}, and turns out to be the conditional
expectation:
\begin{equation}\label{eq:optimal_denoiser}
    \bar{\vx}^{\star}(t, \vxi) \doteq \Ex[\vx \mid \vx_{t} = \vxi].
\end{equation}

This expression justifies the notation \(\bar{\vx}\), which computes
a conditional expectation (i.e., conditional mean or conditional
average). In short, it attempts to remove the noise from the noisy
input, outputting the best possible guess (in expectation and with
respect to the \(\ell^{2}\) distance) of the denoised original random variable.

\begin{example}[Denoising Gaussian Noise from a Mixture of
    Gaussians]\label{example:denoising_gaussian_mixture}
    In this example we compute the Bayes optimal denoiser for an
    important class of distributions, the Gaussian mixture
    model in \Cref{ex:gaussian_mixture}. Let us define \(\phi(\vx ;
    \vmu, \vSigma)\) as the probability density of \(\dNorm(\vmu,
    \vSigma)\) evaluated at \(\vx\). In this notation, the density of
    \(\vx_{t}\) is
    \begin{equation}
        p_{t}(\vx_{t}) = \sum_{k = 1}^{K}\pi_{k}\phi(\vx_{t} ;
        \vmu_{k}, \vSigma_{k} + t^{2}\vI).
    \end{equation}

    Conditioned on \(y\), the variables are jointly Gaussian: if we
    write \(\vx = \vmu_{y} + \vSigma_{y}^{1/2}\vu\) where
    \((\cdot)^{1/2}\) is the matrix square root and \(\vu \sim
    \dNorm(\vzero, \vI)\) independently of \(y\) (and \(\vg\)), then we have
    \begin{equation}
        \mat{\vx \\ \vx_{t}} = \mat{\vmu_{y} \\ \vmu_{y}} +
        \mat{\vSigma_{y}^{1/2} & \vzero \\ \vSigma_{y}^{1/2} &
        t\vI}\mat{\vu \\ \vg}.
    \end{equation}
    This shows that \(\vx\) and \(\vx_{t}\) are jointly Gaussian
    (conditioned on \(y\)) as claimed. Thus we can write
    \begin{equation}
        \mat{\vx \\ \vx_{t}} \sim \dNorm\rp{\mat{\vmu_{y}
            \\ \vmu_{y}}, \mat{\vSigma_{y} & \vSigma_{y} \\ \vSigma_{y} &
        \vSigma_{y} + t^{2}\vI}}.
    \end{equation}
    Thus the conditional expectation of \(\vx\) given \(\vx_{t}\) (i.e., the
    Bayes optimal denoiser conditioned on \(y\)) is famously
    (see also \Cref{exercise:conditional_gaussian})
    \begin{equation}
        \Ex[\vx \mid \vx_{t}, y] = \vmu_{y} + \vSigma_{y}(\vSigma_{y}
        + t^{2}\vI)^{-1}(\vx_{t} - \vmu_{y}).
    \end{equation}
    To find the overall Bayes optimal denoiser, we use the law of
    iterated expectation, obtaining
    \begin{align}
        \bar{\vx}^{\star}(t, \vx_{t})
        &= \Ex[\vx \mid \vx_{t}] \\
        &= \Ex[\Ex[\vx \mid \vx_{t}, y] \mid \vx_{t}] \\
        &= \sum_{k = 1}^{K}\Pr[y = k \mid \vx_{t}]\Ex[\vx \mid \vx_{t}, y = k].
    \end{align}
    The probability can be dealt with as follows. Let \(p_{t \mid
    y}\) be the probability density of \(\vx_{t}\) conditioned on the
    value of \(y\). Then
    \begin{align}
        \Pr[y = k \mid \vx_{t}]
        &= \frac{p_{t \mid y}(\vx_{t} \mid k)\pi_{k}}{p_{t}(\vx_{t})} \\
        &= \frac{\pi_{k}\phi(\vx_{t} ; \vmu_{k}, \vSigma_{k}
        + t^{2}\vI)}{\sum_{i = 1}^{K}\pi_{i}\phi(\vx_{t} ; \vmu_{i},
        \vSigma_{i} + t^{2}\vI)}.
    \end{align}
    On the other hand, the conditional expectation is as described before:
    \begin{equation}
        \Ex[\vx \mid \vx_{t}, y = k] = \vmu_{k} +
        \vSigma_{k}(\vSigma_{k} + t^{2}\vI)^{-1}(\vx_{t} - \vmu_{k}).
    \end{equation}
    So putting this all together, the true Bayes optimal denoiser is
    \begin{align}\label{eq:gmm_bayes_optimal_denoiser}
        \bar{\vx}^{\star}(t, \vx_{t}) &= \sum_{k =
        1}^{K}\Bigg\{\frac{\pi_{k}\phi(\vx_{t}
            ;\vmu_{k}, \vSigma_{k} + t^{2}\vI)}{\sum_{i =
                1}^{K}\pi_{i}\phi(\vx_{t}
            ;\vmu_{i}, \vSigma_{i} + t^{2}\vI)}\\
            &\qquad \qquad \qquad \qquad   \cdot \bp{\vmu_{k} +
        \vSigma_{k}(\vSigma_{k} + t^{2}\vI)^{-1}(\vx_{t} - \vmu_{k})}\Bigg\}.
    \end{align}
    This example is particularly important, and several special cases
    will give us great conceptual insight later. For now, let us
    attempt to extract some geometric intuition from the functional
    form of the optimal denoiser \eqref{eq:gmm_bayes_optimal_denoiser}.

    To understand \eqref{eq:gmm_bayes_optimal_denoiser}
    intuitively, let us first set \(K = 1\) (i.e., one Gaussian) such
    that \(\vx \sim \dNorm(\vmu, \vSigma)\). Let us then diagonalize
    \(\vSigma = \vV\vLambda \vV^{\top}\). Then the Bayes optimal denoiser is
    \begin{align}
        \bar{\vx}^{\star}(t, \vx_{t}) & = \vmu + \vSigma(\vSigma +
        t^{2}\vI)^{-1}(\vx_{t} - \vmu) \\
        & = \vmu + \vV\mat{\lambda_{1}/(\lambda_{1} + t^{2}) & & \\ &
            \ddots & \\ & & \lambda_{D}/(\lambda_{D} +
        t^{2})}\vV^{\top}(\vx_{t} - \vmu),\label{eq:gaussian_denoiser}
    \end{align}
    where \(\lambda_{1}, \dots, \lambda_{D}\) are the eigenvalues of
    \(\vSigma\). We can observe that this denoiser has three steps:
    \begin{itemize}
        \item Translate the input \(\vx_{t}\) by \(\vmu\).
        \item Contract the (translated) input \(\vx_{t} - \vmu\) in
            each eigenvector direction by a quantity
            \(\lambda_{i}/(\lambda_{i} + t^{2})\). If the translated
            input is low-rank and some eigenvalues of \(\vSigma\) are
            zero, these directions get immediately contracted to
            \(0\) by the denoiser, ensuring that the output of the
            contraction is similarly low-rank.
        \item Translate the output back by \(\vmu\).
    \end{itemize}
    This is the geometric interpretation of the denoiser of a
    \textit{single} Gaussian. The overall denoiser of the Gaussian
    mixture model \eqref{eq:gmm_bayes_optimal_denoiser} uses \(K\)
    such denoisers, weighting their output by the posterior
    probabilities \(\Pr[y = k \mid \vx_{t}]\). If the means of the
    Gaussians are well separated, these posterior probabilities are
    very close to \(0\) or \(1\) near each mean or cluster. In this
    regime, the overall denoiser
    \eqref{eq:gmm_bayes_optimal_denoiser} has the same geometric
    interpretation as the above single Gaussian denoiser.
\end{example}

\begin{figure}
    \centering
    \includegraphics[width=\textwidth]{\toplevelprefix/chapters/chapter3/figs/ve_forward_diffusion_denoising.png}\vspace{-0.15in}
    \caption{\small \textbf{Bayes-optimal denoiser and score of a
        Gaussian mixture model.} In the same setting as
        \Cref{fig:ve_forward_density}, we demonstrate the effect of the
        Bayes-optimal denoiser \(\bar{\vx}^{\star}\) by plotting
        \(\vx_{t}\) (red) and \(\bar{\vx}^{\star}(t, \vx_{t})\) (green)
        for some choice of \(t\) and \(\vx_{t}\). By Tweedie's formula
        \Cref{thm:tweedie}, the residual between them is proportional to
        the so-called (Hyv\"arinen) score \(\nabla_{\vx_{t}}\log
        p_{t}(\vx_{t})\). We can see that the score points toward the
    modes of the distribution of \(\vx_{t}\).}
    \label{fig:ve_forward_denoising}
\end{figure}

Intuitively, and as we can see from
\Cref{example:denoising_gaussian_mixture}, the Bayes-optimal denoiser
\(\bar{\vx}^{\star}(t, \cdot)\) should move its input \(\vx_{t}\)
toward the modes of the distribution of \(\vx\). It turns out that we
can quantify this by showing that the Bayes-optimal denoiser
\textit{takes a gradient-ascent step} on the (log-)density
of \(\vx_{t}\), which (recall) we denoted \(p_{t}\). That is,
following the denoiser means moving from the input iterate to a
region of higher probability within this (perturbed) distribution.
For small \(t\), the perturbation is small, so our initial intuition
is therefore almost exactly right. The picture is visualized in
\Cref{fig:ve_forward_denoising} and rigorously formulated as
Tweedie's formula \cite{Robbins1956AnEB}.
\begin{theorem}[Tweedie's Formula]\label{thm:tweedie}
    Suppose that \((\vx_{t})_{t \in [0, T]}\) obeys
    \eqref{eq:additive_gaussian_noise_model}. Let \(p_{t}\) be the
    density of \(\vx_{t}\) (as previously declared). Then
    \begin{equation}\label{eq:tweedie}
        \Ex[\vx \mid \vx_{t}] = \vx_{t} + t^{2}\nabla_{\vx_{t}} \log
        p_{t}(\vx_{t}).
    \end{equation}
\end{theorem}
\begin{proof}
    For the proof, suppose that \(\vx\) has a
    density,\footnote{Note that the
        theorem is true without this assumption, and it can be proven through
        a minimal modification of this proof to use probability measures
    instead of densities.} and call this density \(p\). Let
    \(p_{0 \mid t}\) and \(p_{t \mid 0}\) be the conditional densities of \(\vx
    = \vx_{0}\) given \(\vx_{t}\) and \(\vx_{t}\) given \(\vx\) respectively.
    Let \(\phi(\vx ; \vmu, \vSigma)\) be the density of \(\dNorm(\vmu,
    \vSigma)\) evaluated at \(\vx\), so that \(p_{t \mid 0}(\vx_{t} \mid \vx)
    = \phi(\vx_{t} ; \vx, t^{2}\vI)\). Then a simple calculation gives
    \begin{align}
        \nabla_{\vx_{t}} \log p_{t}(\vx_{t})
        &= \frac{\nabla_{\vx_{t}}p_{t}(\vx_{t})}{p_{t}(\vx_{t})} \\
        &=
        \frac{1}{p_{t}(\vx_{t})}\nabla_{\vx_{t}}\int_{\R^{D}}p(\vx)p_{t
        \mid 0}(\vx_{t} \mid \vx)\odif{\vx} \\
        &=
        \frac{1}{p_{t}(\vx_{t})}\nabla_{\vx_{t}}\int_{\R^{D}}p(\vx)\phi(\vx_{t}
        ;\vx, t^{2}\vI)\odif{\vx} \\
        &=
        \frac{1}{p_{t}(\vx_{t})}\int_{\R^{D}}p(\vx)[\nabla_{\vx_{t}}\phi(\vx_{t}
        ;\vx, t^{2}\vI)]\odif{\vx} \\
        &= \frac{1}{p_{t}(\vx_{t})}\int_{\R^{D}}p(\vx)\phi(\vx_{t} ;
        \vx, t^{2}\vI)\bs{-\frac{\vx_{t} - \vx}{t^{2}}}\odif{\vx} \\
        &=
        \frac{1}{t^{2}p_{t}(\vx_{t})}\int_{\R^{D}}p(\vx)\phi(\vx_{t}
        ; \vx, t^{2}\vI)[\vx - \vx_{t}]\odif{\vx} \\
        &= \frac{1}{t^{2}p_{t}(\vx_{t})}\int_{\R^{D}}p(\vx)\phi(\vx_{t} ; \vx,
        t^{2}\vI)\vx \odif{\vx} \\
        & \quad -
        \frac{\vx_{t}}{t^{2}p_{t}(\vx_{t})}\int_{\R^{D}}p(\vx)\phi(\vx_{t}
        ; \vx, t^{2}\vI)\odif{\vx}\nonumber \\
        &= \frac{1}{t^{2}p_{t}(\vx_{t})}\int_{\R^{D}}p(\vx)p_{t \mid
        0}(\vx_{t} \mid \vx)\vx \odif{\vx} -
        \frac{\vx_{t}}{t^{2}p_{t}(\vx_{t})}p_{t}(\vx_{t}) \\
        &=
        \frac{1}{t^{2}p_{t}(\vx_{t})}\int_{\R^{D}}p_{t}(\vx_{t})p_{0
        \mid t}(\vx \mid \vx_{t})\vx \odif{\vx} -
        \frac{\vx_{t}}{t^{2}p_{t}(\vx_{t})}p_{t}(\vx_{t}) \\
        &= \frac{1}{t^{2}}\int_{\R^{D}}p_{0 \mid t}(\vx \mid
        \vx_{t})\vx \odif{\vx} - \frac{\vx_{t}}{t^{2}} \\
        &= \frac{1}{t^{2}}\Ex[\vx \mid \vx_{t}] - \frac{\vx_{t}}{t^{2}} \\
        &= \frac{\Ex[\vx \mid \vx_{t}] - \vx_{t}}{t^{2}}.
    \end{align}
    Simple rearrangement of the above equality proves the theorem.
\end{proof}
This result develops a connection between denoising and optimization:
the Bayes-optimal denoiser takes a single step of gradient ascent on
the perturbed data density \(p_{t}\), and the step size adaptively
becomes smaller (i.e., takes more precise steps) as the perturbation
to the data distribution grows smaller. The quantity
\(\nabla_{\vx_{t}}\log p_{t}(\vx_{t})\) is called the
\textit{(Hyv\"arinen) score} and frequently appears in discussions
about denoising, etc.; it first appeared in a paper of Aapo
Hyv\"arinen in the context of ICA \cite{hyvarinen05a}.

Similar to how one step of gradient descent is almost never sufficient to
minimize an objective in practice when initializing far from the optimum, the
output of the Bayes-optimal denoiser \(\bar{\vx}^{\star}(t, \cdot)\)
is almost never contained in a high-probability region of the data
distribution when \(t\) is large, \textit{especially} when the data
have low-dimensional structure. We illustrate this point explicitly
in the following example.
\begin{example}[Denoising a two-point
    mixture]\label{example:denoising_twopoints}
    Let \(x\) be uniform on the two-point set \(\{-1, +1\}\) and let
    \((x_{t})_{t \in [0, T]}\) follow
    \eqref{eq:additive_gaussian_noise_model}. This is precisely a
    degenerate Gaussian mixture model with priors equal to
    \(\frac{1}{2}\), means \(\{-1, +1\}\), and covariances both equal
    to \(0\). For a fixed \(t > 0\) we can use the calculation of the
    Bayes-optimal denoiser in \eqref{eq:gmm_bayes_optimal_denoiser}
    to obtain (proof as exercise)
    \begin{equation}
        \bar{x}^{\star}(t, x_{t}) = \frac{\phi(x_{t} ; +1, t^{2}) - \phi(x_{t}
        ; -1, t^{2})}{\phi(x_{t} ; 1, t^{2}) + \phi(x_{t} ; -1,
        t^{2})} = \tanh\rp{\frac{x_{t}}{t^{2}}}.
    \end{equation}
    For \(t\) near \(0\), this quantity is near \(\{-1, +1\}\) for
    almost all inputs \(\bar{x}^{\star}(t, x_{t})\). However, for
    \(t\) large, this quantity is not necessarily even approximately
    in the original support of \(x\), which, remember, is \(\{-1,
    +1\}\). In particular, for \(x_{t} \approx 0\) it holds
    \(\bar{x}^{\star}(t, x_{t}) \approx 0\), which lies completely in
    between the two possible points. Thus \(\bar{x}^{\star}\)
    \textit{will not output ``realistic'' \(x\)}. More rigorously,
    the distribution of \(\bar{x}^{\star}(t, x_{t})\) is very
    different from the distribution of \(x\).
\end{example}

\subsubsection{Incremental Denoising}
Therefore, to denoise the very noisy sample \(\vx_{T}\) (where,
recall, \(T\) is the maximum time), we cannot use the denoiser just
once. Instead, we must apply it many times, analogously to gradient
descent with decaying step sizes, to converge to a stationary point
\(\hat{\vx}\). Namely, we use the denoiser to go from \(\vx_{T}\) to
\(\hat{\vx}_{T - \delta}\), which approximates \(\vx_{T - \delta}\),
then from \(\hat{\vx}_{T - \delta}\) to \(\hat{\vx}_{T - 2\delta}\),
and so on, all the way from \(\hat{\vx}_{\delta}\) to \(\hat{\vx} =
\hat{\vx}_{0}\). Each denoising step makes the action of the denoiser
more like a gradient step on the original (log-)density.

More formally, we uniformly discretize \([0, T]\) into \(L + 1\)
timesteps \(0 = t_{0} < t_{1} < \cdots < t_{L} = T\), i.e.,
\begin{equation}
    t_{\ell} = \frac{\ell}{L}T, \qquad \ell \in \{0, 1, \dots, L\}.
\end{equation}
Then for each \(\ell \in [L] = \{1, 2, \dots, L\}\), going from
\(\ell = L\) to \(\ell = 1\), we run the iteration
\begin{align}
    \hat{\vx}_{t_{\ell - 1}}
    &= \Ex[\vx_{t_{\ell - 1}} \mid \vx_{t_{\ell}} = \hat{\vx}_{t_{\ell}}] \\
    &= \Ex[\vx + t_{\ell - 1}\vg \mid \vx_{t_{\ell}} = \hat{\vx}_{t_{\ell}}] \\
    &= \Ex\rs{\vx + t_{\ell - 1}\cdot\frac{\vx_{t_{\ell}} -
    \vx}{t_{\ell}} \given \vx_{t_{\ell}} = \hat{\vx}_{t_{\ell}}} \\
    &= \frac{t_{\ell - 1}}{t_{\ell}}\hat{\vx}_{t_{\ell}} + \bp{1 -
    \frac{t_{\ell - 1}}{t_{\ell}}}\Ex[\vx \mid \vx_{t_{\ell}} =
    \hat{\vx}_{t_{\ell}}] \\
    &= \bp{1 - \frac{1}{\ell}} \hat{\vx}_{t_{\ell}} + \frac{1}{\ell}
    \bar{\vx}^{\star}(t_{\ell}, \hat{\vx}_{t_{\ell}}). \label{eq:iter denoise}
\end{align}

At the beginning of the iteration, when \(\ell\) is large, we barely
trust the denoiser's output and mostly keep the current iterate.
This makes sense, as the denoiser can have huge variance (cf.
\Cref{example:denoising_twopoints}). When \(\ell\) is small, the
denoiser ``locks on'' to the modes of the data distribution, since a
denoising step is essentially a gradient step on the true
distribution's log-density, and we can trust it not to produce
unreasonable samples; thus the denoising step is dominated by the
denoiser's output. At \(\ell = 1\) we even discard the
current iterate and keep only the denoiser's output.

\begin{figure}[t]
    \centering
    \includegraphics[width=\textwidth]{\toplevelprefix/chapters/chapter3/figs/ve_gmm_denoising.png}
    \caption{\small\textbf{Denoising a low-rank mixture of
        Gaussians.} Each figure shows samples from the true data
        distribution (gray, orange, red) and samples undergoing the
        denoising process \eqref{eq:denoising-iteration-basic} (light
        blue). At top left, the process has just started and the noise
        is very large. As the process continues, the noise is pushed
        further toward the support of the low-rank data distribution.
        Finally, in the bottom right, the generated samples are perfectly
        aligned with the support of the data and closely resemble
    samples drawn from the low-rank Gaussian mixture model.}
    \label{fig:ve_gmm_denoising}
\end{figure}

We visualize the convergence process in \(\R^{3}\) in
\Cref{fig:ve_gmm_denoising}. While we will develop some rigorous
results about convergence later, we now give a closed-form example of
the iterative denoising process applied to denoising a Gaussian data
distribution.

\begin{example}[Iterative denoising for a single Gaussian]
    \label{ex:single_gaussian_denoising}
    To understand the above iterative denoising iteration
    intuitively, we apply it to denoise \(\vx_{T}\) defined in
    \eqref{eq:additive_gaussian_noise_model}, where \(\vx \sim
    \dNorm(\vmu, \vSigma)\). Using \eqref{eq:iter denoise}, we compute
    \begin{align}
        \hat{\vx}_{t_{\ell - 1}} - \vmu
        &= \bp{1 - \frac{1}{\ell}}\hat{\vx}_{t_{\ell}} +
        \frac{1}{\ell}\bar{\vx}^{\star}(t_{\ell},
        \hat{\vx}_{t_{\ell}}) - \vmu \\
        &= \bp{1 - \frac{1}{\ell}}(\hat{\vx}_{t_{\ell}} - \vmu) +
        \frac{1}{\ell}(\bar{\vx}^{\star}(t_{\ell},
        \hat{\vx}_{t_{\ell}}) - \vmu).
    \end{align}
    Then, plugging in the form for \(\bar{\vx}^{\star}\) that we
    obtain from \eqref{eq:gaussian_denoiser} and performing some
    tedious matrix algebra, we obtain
    \begin{equation}
        \vx_{t_{\ell - 1}} - \vmu = \vV\mat{1 - \frac{\ell
            T^{2}}{\lambda_{1}L^{2} + \ell^{2} T^{2}} & & \\ & \ddots &
            \\ & & 1 - \frac{\ell T^{2}}{\lambda_{D}L^{2} +
        \ell^{2}T^{2}}}\vV^{\top}(\vx_{t_{\ell}} - \vmu).
    \end{equation}
    Notice that all entries in this diagonal matrix are in \((0,
    1)\). Taking the \(\ell^{2}\)-norm and using a cheap yet tight
    upper bound, it holds
    \begin{equation}\label{eq:denoising_contraction_1step}
        \norm{\vx_{t_{\ell - 1}} - \vmu}_{2} \leq \underbrace{\max_{i
            \in [D]}\bc{1 - \frac{\ell T^{2}}{\lambda_{i}L^{2} +
        \ell^{2}T^{2}}}}_{c_{\ell, L}} \cdot \norm{\vx_{t_{\ell}} - \vmu}_{2}.
    \end{equation}
    At first glance, this inequality seems to imply that the
    denoising iteration is \textit{very close}, conceptually, to a
    \textit{contraction mapping} such as power iteration
    \eqref{eq:power_iteration} sending the iterate \(\vx_{t}\) to
    the mean \(\vmu\), since \(c_{\ell, L}\) are all in \((0, 1)\)
    for any \(\ell\) and \(L\). However, this is false. To see why,
    observe that by iterating \eqref{eq:denoising_contraction_1step} it holds
    \begin{equation}
        \norm{\vx_{0} - \vmu}_{2} \leq \prod_{\ell = 1}^{L}c_{\ell,
        L} \cdot \norm{\vx_{T} - \vmu}_{2}.
    \end{equation}
    If each denoising update had similar properties to a contraction
    mapping, then \(\prod_{\ell = 1}^{L}c_{\ell, L}\) would go to
    \(0\) as \(L \to \infty\), and the iterate would converge to
    \(\vmu\) extremely fast, rendering the sampler completely useless
    or vacuous. But in reality the denoising iteration is
    \textit{actually not a contraction mapping} since \(c_{\ell, L}\)
    changes with both the indices \(\ell\) and \(L\), so that \(\prod_{\ell =
    1}^{L}c_{\ell, L}\) does not go to \(0\) as \(L \to \infty\), as
    demonstrated in \Cref{fig:not_contraction_mapping} (it is also
    possible to compute the limit analytically). In reality, the
    sampler ensures that the terminal iterates spread out amongst the
    support of the data distribution, enabling actual sampling from
    the data distribution, which is quite unlike the objective of
    power iteration (for example).
    \begin{figure}
        \centering
        \includegraphics[width=\textwidth]{\toplevelprefix/chapters/chapter3/figs/not_contraction_mapping.png}
        \caption{\small\textbf{The denoising iteration from
                \Cref{ex:single_gaussian_denoising} is not a contraction
            mapping.} Plot of the contraction coefficient products
            \(\prod_{\ell = 1}^{L}c_{\ell, L}\) in the case where the
            data are in \(D = 1\), the time horizon \(T = 1\), and the
            data variance \(\lambda_{1} = a\), for varying values of
            \(a\). If the denoising iteration were (conceptually similar
            to) a contraction mapping, these products would go to \(0\),
        but instead they converge to non-zero values.}
        \label{fig:not_contraction_mapping}
    \end{figure}
\end{example}

Recall that we wanted to build the iterative denoising process to
reduce the entropy of the iterates. While we did this in a roundabout
way by inverting a process that adds entropy, it is now time to pay the
piper and confirm that our iterative denoising process reduces the
entropy.

\begin{theorem}[Simplified Version of \Cref{thm:conditioning_reduces_entropy}]
    Suppose that \((\vx_{t})_{t \in [0, T]}\) obeys
    \eqref{eq:additive_gaussian_noise_model}. Then, under certain
    technical conditions on \(\vx\), for every \(s < t\) with \(s, t
    \in (0, T]\),
    \begin{equation}
        h(\Ex[\vx_{s} \mid \vx_{t}]) < h(\vx_{t}).
    \end{equation}
\end{theorem}
The full statement of the theorem and the proof itself require
some technicality, so they are postponed to
\Cref{sub:denoising_entropy_decreases}.

\subsubsection{Computing the Denoiser}
The last thing we discuss here is that many times, we will
\textit{not be able to compute} \(\bar{\vx}^{\star}(t, \cdot)\) for
any \(t\), since we do not have the distribution \(p_{t}\). But we
can try to \textit{learn one from data}. Recall that the denoiser
\(\bar{\vx}^{\star}\) is defined in \eqref{eq:denoising_loss} as
minimizing the mean-squared error \(\Ex\norm{\bar{\vx}(t, \vx_{t}) -
\vx}_{2}^{2}\). We can use this mean-squared error as a loss or
objective function to learn the denoiser. For example, we can
parameterize \(\bar{\vx}(t, \cdot)\) by a neural network, writing it
as \(\bar{\vx}_{\theta}(t, \cdot)\), and optimize the loss over the
parameter space \(\Theta\):
\begin{equation}
    \min_{\theta \in \Theta}\Ex\norm{\bar{\vx}_{\theta}(t, \vx_{t}) -
    \vx}_{2}^{2}.
\end{equation}
The solution to this optimization problem, implemented via gradient
descent or a similar algorithm, will give us a
\(\bar{\vx}_{\theta^{\star}}(t, \cdot)\) which is a good approximation
to \(\bar{\vx}^{\star}(t, \cdot)\) (at least if the training works)
and which we will use as our denoiser.

What is a good architecture for this neural network
\(\bar{\vx}_{\theta^{\star}}(t, \cdot)\)? To answer this question, we
will first attempt to understand the case of data \(\vx\) that are
\textit{low-rank}, i.e., having low-rank support. The most canonical
way to represent low-rank data is through a \textit{low-rank
Gaussian}, and so we start our exploration by assuming that \(\vx\)
is drawn (approximately) from a low-rank Gaussian:
\begin{equation}\label{eq:x_gaussian_distribution}
    \vx \sim \dNorm(\vzero, \vU\vU^{\top}),
\end{equation}
where \(\vU \in \O(D, P)\) is a tall orthogonal matrix with \(P \leq
D\) (using zero-mean for essentially notational convenience).  The
optimal denoiser under \eqref{eq:x_gaussian_distribution} is (from
\Cref{example:denoising_gaussian_mixture}):
\begin{equation}\label{eq:optimal denoiser}
    \bar{\vx}^{\star}(t, \vx_{t}) = \vU\vU^{\top}(\vU\vU^{\top} +
    t^{2}\vI)^{-1}\vx_{t}.
\end{equation}
The matrix  \(\vU\vU^{\top} + t^{2}\vI\) is a low-rank perturbation
of the full-rank matrix \(t^{2}\vI\), and thus ripe for
simplification via the \textit{Sherman-Morrison-Woodbury identity},
i.e., for matrices \(\vA, \vC, \vU, \vV\) such that \(\vA\) and
\(\vC\) are invertible,
\begin{equation}\label{eq:sherman_morrison_woodbury}
    (\vA + \vU\vC\vV)^{-1} = \vA^{-1} - \vA^{-1}\vU(\vC^{-1} +
    \vV\vA^{-1}\vU)^{-1}\vV\vA^{-1}.
\end{equation}
We prove this identity in
\Cref{exercise:sherman_morrison_woodbury_identity}. For now we apply
this identity with \(\vA = t^{2}\vI\), \(\vU = \vU\), \(\vV =
\vU^{\top}\), and \(\vC = \vI\), obtaining
\begin{align}
    (\vU\vU^{\top} + t^{2}\vI)^{-1}
    &= \frac{1}{t^{2}}\vI - \frac{1}{t^{4}}\vU\bp{\vI +
    \frac{1}{t^{2}}\vU^{\top}\vU}^{-1}\vU^{\top} \\
    &= \frac{1}{t^{2}}\vI - \frac{1}{t^{4}\bp{1 +
    \frac{1}{t^{2}}}}\vU\vU^{\top} \\
    &= \frac{1}{t^{2}}\bp{\vI - \frac{1}{1 +
    t^{2}}\vU\vU^{\top}}.\label{eq:gaussian_lowrank_inverse_perturbation_simplified}
\end{align}
This implies
\begin{align}
    \vU\vU^{\top}(\vU\vU^{\top} + t^{2}\vI)^{-1}
    &= \frac{1}{t^{2}}\vU\vU^{\top}\bp{\vI - \frac{1}{1 +
    t^{2}}\vU\vU^{\top}} \\
    &= \frac{1}{t^{2}}\bp{1 - \frac{1}{1 + t^{2}}}\vU\vU^{\top} \\
    &= \frac{1}{1 +
    t^{2}}\vU\vU^{\top}.\label{eq:gaussian_lowrank_denoiser_simplified}
\end{align}
Our optimal denoiser is a (rescaled) linear projection onto the
low-rank subspace supporting our data distribution. This is
essentially the same thing as probabilistic PCA from
\Cref{ch:classic-models} (cf.~\Cref{subsec:probabilistic PCA}). We can
formalize this intuition by considering a class of
subspace-parameterized denoisers:
\begin{equation}\label{eq:optimal nn}
    \bar{\vx}_{\vV}(t, \vx_{t}) = \frac{1}{1 + t^{2}}\vV\vV^{\top}\vx_{t},
\end{equation}
and showing that the denoising problem \(\min_{\vV \in \O(D,
    P)}\Ex\norm{\bar{\vx}_{\vV}(t, \vx_{t}) -
\vx}_{2}^{2}\) obtains the same solution as probabilistic PCA
(\Cref{subsec:probabilistic PCA}). Indeed, plugging our subspace
denoiser into the training loss \eqref{eq:denoising_loss}, we obtain
\begin{equation}
    \min_{\vV \in \O(D, P)}\Ex\norm{\bar{\vx}_{\vV}(t,
    \vx_{t}) - \vx}_{2}^{2} = \Ex\norm*{\frac{1}{1 +
    t^{2}}\vV\vV^{\top}(\vx + t\vg) - \vx}_{2}^{2},
\end{equation}
where the equality is due to
\eqref{eq:additive_gaussian_noise_model}. Conditioned on \(\vx\)
(thus integrating over \(\vg\)), we compute
\begin{align}
    & \Ex_{\vg}\norm*{\frac{1}{1 + t^{2}}\vV\vV^{\top}(\vx + t\vg) -
    \vx}_{2}^{2} \\
    &= \norm*{\frac{1}{1 + t^{2}}\vV\vV^{\top}\vx - \vx}_{2}^{2} -
    \frac{t}{1+t^2}\Ex_{\vg} \ip*{\vx - \frac{1}{1 +
    t^{2}}\vV\vV^{\top}\vx}{\vV\vV^{\top}\vg} \\
    &\quad + \frac{t^2}{(1+t^2)^2} \Ex_{\vg}\norm*{\vV\vV^{\top} \vg}_{2}^{2} \\
    &= \norm*{\frac{1}{1 + t^{2}}\vV\vV^{\top}\vx - \vx}_{2}^{2} +
    \frac{t^2P}{(1+t^2)^2}
\end{align}
where the second equality follows from \(\vg \sim \dNorm(\vzero,
\vI)\) and the fact that \(\vV \in \O(D, P)\) so that \(
    \Ex_{\vg}\norm{\vV\vV^\top
    \vg}_{2}^{2} =
\Ex_{\vg}\norm{\vV^\top \vg}_{2}^{2} = P\). Therefore, the denoising
problem is equivalent to
\begin{align}
    \min_{\vV \in \O(D, P)} & \Bigg\{\Ex \norm*{\frac{1}{1 +
        t^{2}}\vV\vV^{\top}\vx - \vx}_{2}^{2} \\
        & \qquad = \Ex\norm{\vx}_{2}^{2} + \bp{\frac{1}{(1 +
            t^{2})^{2}} - \frac{2}{1 +
    t^{2}}}\Ex\norm{\vV^{\top}\vx}_{2}^{2}\Bigg\},
\end{align}
which in turn is equivalent to the problem
\begin{equation}\label{eq:PCA}
    \max_{\vV \in \O(D, P)}\Ex\norm{\vV^{\top}\vx}_{2}^{2}
\end{equation}
which was shown in \Cref{sub:pca} to be equivalent to the problem
\begin{equation}
    \min_{\vV \in \O(D, P)}\Ex\norm{\vx - \vV\vV^{\top}\vx}_{2}^{2},
\end{equation}
i.e., the probabilistic PCA formulation \eqref{eq:PPCA}.

Thus, if we were to assume our data were drawn from a low-rank
Gaussian, our denoising algorithm would essentially consist of many
PCA iterations in a row applied to progressively less-noisy data.
While this method sounds very simple, it can actually produce complex
and relatively powerful denoising behaviors when applied to practical
data, as demonstrated by PCANet \parencite{chan2015pcanet}.
Nevertheless, to obtain a more general denoising operator
that may work on complex and multi-modal data distributions, we will
consider a broader class of data distributions: the ubiquitous
class of \textit{low-rank Gaussian mixture models}, whose denoiser
we computed in \Cref{ex:gaussian_mixture} and
\Cref{example:denoising_gaussian_mixture}. In the context of
denoising, since Gaussian mixture models can approximate a very broad
class of distributions arbitrarily well, in practice optimizing among
the class of denoisers for Gaussian mixture models can potentially
give us something close to the optimal denoiser for the real data distribution.

Formally, we now assume that \(\vx\) is drawn from a low-rank
Gaussian mixture model \(\frac{1}{K}\sum_{k = 1}^{K}\dNorm(\vzero,
\vU_{k}\vU_{k}^{\top})\), where \(\vU_{k} \in \O(D, P)\) are tall
orthogonal matrices whose columns span the \(P\)-dimensional
subspaces that support the \(K\) low-rank Gaussian components. We write
\begin{equation}\label{eq:MoG1}
    \vx \sim \frac{1}{K}\sum_{k = 1}^{K}\dNorm(\vzero, \vU_{k}\vU_{k}^{\top}).
\end{equation}
Then the optimal denoiser under
\eqref{eq:additive_gaussian_noise_model} is (from
\Cref{example:denoising_gaussian_mixture})
\begin{align}
    \bar{\vx}^{\star}(t, \vx_{t})
    &= \sum_{k = 1}^{K}\Bigg\{\frac{\phi(\vx_{t} ; \vzero,
        \vU_{k}\vU_{k}^{\top} + t^{2}\vI)}{\sum_{i = 1}^{K}\phi(\vx_{t} ;
                \vzero, \vU_{i}\vU_{i}^{\top} +
        t^{2}\vI)} \\
        &\qquad \qquad \qquad \qquad
        \cdot\bp{\vU_{k}\vU_{k}^{\top}(\vU_{k}\vU_{k}^{\top} +
    t^{2}\vI)^{-1}\vx_{t}}\Bigg\}.
\end{align}
Applying the same Sherman-Morrison-Woodbury identity as in
\eqref{eq:gaussian_lowrank_inverse_perturbation_simplified}, we obtain
\begin{equation}
    (\vU_{k}\vU_{k}^{\top} + t^{2}\vI)^{-1} = \frac{1}{t^{2}}\bp{\vI
    - \frac{1}{1 + t^{2}}\vU_{k}\vU_{k}^{\top}}.
\end{equation}
We can then compute the posterior probabilities as follows. Since the
\(\vU_{k}\) are orthogonal, \(\det(\vU_{k}\vU_{k}^{\top}
+ t^{2}\vI)\) is the same for every \(k\). Thus
\begin{align}
    &\frac{\phi(\vx_{t} ; \vzero, \vU_{k}\vU_{k}^{\top} + t^{2}\vI)}{\sum_{i
    = 1}^{K}\phi(\vx_{t} ; \vzero, \vU_{i}\vU_{i}^{\top} + t^{2}\vI)} \\
    &= \frac{\exp\rp{-\frac{1}{2}\vx_{t}^{\top}(\vU_{k}\vU_{k}^{\top}
    + t^{2}\vI)^{-1}\vx_{t}}}{\sum_{i =
        1}^{K}\exp\rp{-\frac{1}{2}\vx_{t}^{\top}(\vU_{i}\vU_{i}^{\top} +
    t^{2}\vI)^{-1}\vx_{t}}} \\
    &= \frac{\exp\rp{-\frac{1}{2t^{2}}\vx_{t}^{\top}\bp{\vI -
    \frac{1}{1 + t^{2}}\vU_{k}\vU_{k}^{\top}}\vx_{t}}}{\sum_{i =
        1}^{K}\exp\rp{-\frac{1}{2t^{2}}\vx_{t}^{\top}\bp{\vI - \frac{1}{1
    + t^{2}}\vU_{i}\vU_{i}^{\top}}\vx_{t}}} \\
    &= \frac{\exp\rp{-\frac{1}{2t^{2}}\norm{\vx_{t}}_{2}^{2} +
            \frac{1}{2t^{2}(1 +
    t^{2})}\norm{\vU_{k}^{\top}\vx_{t}}_{2}^{2}}}{\sum_{i =
        1}^{K}\exp\rp{-\frac{1}{2t^{2}}\norm{\vx_{t}}_{2}^{2} +
    \frac{1}{2t^{2}(1 + t^{2})}\norm{\vU_{i}^{\top}\vx_{t}}_{2}^{2}}} \\
    &=
    \frac{\exp\rp{-\frac{1}{2t^{2}}\norm{\vx_{t}}_{2}^{2}}\exp\rp{\frac{1}{2t^{2}(1
                    +
    t^{2})}\norm{\vU_{k}^{\top}\vx_{t}}_{2}^{2}}}{\exp\rp{-\frac{1}{2t^{2}}\norm{\vx_{t}}_{2}^{2}}\sum_{i
        = 1}^{K}\exp\rp{\frac{1}{2t^{2}(1 +
    t^{2})}\norm{\vU_{i}^{\top}\vx_{t}}_{2}^{2}}} \\
    &= \frac{\exp\rp{\frac{1}{2t^{2}(1 +
    t^{2})}\norm{\vU_{k}^{\top}\vx_{t}}_{2}^{2}}}{\sum_{i =
        1}^{K}\exp\rp{\frac{1}{2t^{2}(1 +
    t^{2})}\norm{\vU_{i}^{\top}\vx_{t}}_{2}^{2}}}.
\end{align}
This is a softmax operation weighted by the projection of \(\vx_{t}\)
onto each subspace measured by \(\norm{\vU_{i}^{\top}\vx_{t}}_{2}\)
(tempered by the temperature \(2t^{2}(1 + t^{2})\)). Meanwhile, the
component denoisers can be written in the same way as
\eqref{eq:gaussian_lowrank_denoiser_simplified} to obtain
\begin{equation}
    \vU_{k}\vU_{k}^{\top}(\vU_{k}\vU_{k}^{\top} +
    t^{2}\vI)^{-1}\vx_{t} = \frac{1}{1 + t^{2}}\vU_{k}\vU_{k}^{\top}\vx_{t}.
\end{equation}
Putting these together, we have
\begin{equation}\label{eq:gmm_lowrank_denoiser}
    \bar{\vx}^{\star}(t, \vx_{t}) = \frac{1}{1 + t^{2}}\sum_{k =
    1}^{K}\frac{\exp\rp{\frac{1}{2t^{2}(1 +
    t^{2})}\norm{\vU_{k}^{\top}\vx_{t}}_{2}^{2}}}{\sum_{i =
        1}^{K}\exp\rp{\frac{1}{2t^{2}(1 +
    t^{2})}\norm{\vU_{i}^{\top}\vx_{t}}_{2}^{2}}}\vU_{k}\vU_{k}^{\top}\vx_{t},
\end{equation}
i.e., a projection of \(\vx_{t}\) onto each of the \(K\) subspaces,
weighted by a softmax of a quadratic function of \(\vx_{t}\). This
functional form is similar to an \textit{attention
mechanism} in a transformer architecture! As we will see in
\Cref{ch:deep-networks}, this is no coincidence at all; the deep
link between denoising and lossy compression (to be covered in
\Cref{sec:lossy_compression}) makes transformer denoisers so
effective in practice. Thus, our Gaussian mixture model
theory motivates the use of transformer-like neural networks for denoising.

Overall, the learned denoiser architecture corresponds to an
(implicit parametric) encoding scheme of the given data, since it can
be used to denoise/project onto the data distribution. Training a
denoiser is equivalent to finding a better coding scheme, and this
partially implements approach A2 (i.e., pursuing a
distribution by searching for schemes with better coding rates) at
the end of \Cref{sub:min_entropy}. In the sequel, we discuss how
to implement the other approach.

\subsection{Sampling a Distribution via Iterative
Denoising}\label{sub:sampling_denoising}

Remember that at the end of \Cref{sub:min_entropy}, we discussed two
approaches for pursuing a distribution with low-dimensional structure.
The first approach (A1) starts with a normal distribution of high
entropy and gradually reduces its entropy until it reaches the data
distribution. We call this procedure \textit{sampling}, since we
generate new samples. It is now time to discuss how to implement this
with the toolkit we have built.

We know how to denoise very noisy samples \(\vx_{T}\) to obtain
approximations \(\hat{\vx}\) whose distributions resemble that of the
target random variable \(\vx\). Yet the sampling procedure requires us
to start from a template distribution that has \textit{no} influence
from the distribution of \(\vx\) and to use the denoiser to guide the
iterates toward the distribution of \(\vx\). How can we do this? One
way is motivated as follows:
\begin{equation}\label{eq:ve_converge_to_large_gaussian}
    \frac{\vx_{T}}{T} = \frac{\vx + T\vg}{T} = \frac{\vx}{T} + \vg
    \xrightarrow{T\to +\infty} \vg \sim \dNorm(\vzero, \vI).
\end{equation}
Thus, \(\vx_{T} \approx \dNorm(\vzero, T^{2}\vI)\). This
approximation is quite good for almost all practical distributions,
and is visualized in \Cref{fig:xT_vs_noise}.
\begin{figure}
    \centering
    \includegraphics[width=0.75\textwidth]{\toplevelprefix/chapters/chapter3/figs/xT_vs_noise.png}
    \caption{\small\textbf{Visualizing \(\vx_{T}\) versus
        \(\dNorm(\vzero, T^{2}\vI)\).} \textit{Left:} A plot of Gaussian
        mixture model data \(\vx\). \textit{Right:} A plot of \(\vx\) as
        well as \(\vx_{T}\) and an independent sample of \(\dNorm(\vzero,
        T^{2}\vI)\), for \(T = 10\). On the right plot, \(\vx\) is
        plotted in the same colors as the left; however, samples from
        \(\vx_{T}\) and \(\dNorm(\vzero, T^{2}\vI)\) are both much
        larger, on average, than samples from \(\vx\), and so it appears
        much smaller because of the scaling. Despite this large blow-up,
        we clearly observe the similarities in the distributions of
    \(\vx_{T}\) and \(\dNorm(\vzero, T^{2}\vI)\).}
    \label{fig:xT_vs_noise}
\end{figure}

So, discretizing \([0, T]\) into \(0 = t_{0} < t_{1} < \cdots < t_{L}
= T\) uniformly using \(t_{\ell} = T\ell / L\) (as in the previous
section), one possible way to sample from pure noise is:
\begin{itemize}
    \item Sample \(\hat{\vx}_{T} \sim \dNorm(\vzero, T^{2}\vI)\)
        (independently of everything else).
    \item Run the denoising iteration as in
        \Cref{sub:intro_diffusion_denoising}, i.e.,
        \begin{equation}\label{eq:denoising-iteration-basic}
            \hat{\vx}_{t_{\ell - 1}} = \bp{1 -
            \frac{1}{\ell}}\hat{\vx}_{t_{\ell}} +
            \frac{1}{\ell}\bar{\vx}^{\star}(t_{\ell}, \hat{\vx}_{t_{\ell}}).
        \end{equation}
    \item Output \(\hat{\vx} = \hat{\vx}_{0}\).
\end{itemize}
This is conceptually all there is behind \textit{diffusion models},
which transform noise into data samples in accordance with the first
desideratum. However, a few steps remain before we obtain models that
can actually sample from real data distributions like images under
practical resource constraints. In the sequel, we introduce and
motivate several such steps.

\paragraph{Step 1: Different discretizations.} The first step is motivated
by the observation that \textit{we do not need to spend so many
denoising iterations} at large \(t\). \Cref{fig:ve_gmm_denoising}
shows that the first \(200\) or \(300\) iterations out of the \(500\)
iterations of the sampling process are spent contracting the noise
toward the data distribution as a whole, before the remaining
iterations push the samples toward a subspace. Given a fixed
iteration count \(L\), this suggests that we should allocate more
timesteps \(t_{\ell}\) near \(t = 0\) than near \(t = T\). During
sampling (and training), we can therefore use an alternative
discretization of \([0, T]\) into \(0 \leq t_{0} < t_{1} < \cdots <
t_{L} \leq T\), such as an \textit{exponential discretization}:
\begin{equation}\label{eq:denoising_exponential_discretization}
    t_{\ell} = C_{1}(e^{C_{2}\ell} - 1), \qquad \forall \ell \in \{0,
    1, \dots, L\}
\end{equation}
where \(C_{1}, C_{2} > 0\) are constants that can be tuned for
optimal performance in practice; theoretical analyses often specify
such optimal constants as well. The denoising/sampling iteration then
becomes
\begin{equation}\label{eq:denoising_iteration}
    \hat{\vx}_{t_{\ell - 1}} \doteq \frac{t_{\ell -
    1}}{t_{\ell}}\hat{\vx}_{t_{\ell}} + \bp{1 - \frac{t_{\ell -
    1}}{t_{\ell}}}\bar{\vx}^{\star}(t_{\ell}, \hat{\vx}_{t_{\ell}}),
\end{equation}
with \(\hat{\vx}_{t_{L}} \sim \dNorm(\vzero, t_{L}^{2}\vI)\).

\begin{figure}[t]
    \centering
    \includegraphics[width=\textwidth]{\toplevelprefix/chapters/chapter3/figs/vp_gmm_denoising.png}
    \caption{\small\textbf{Denoising a mixture of Gaussians using the
        VP diffusion process.} We use the same figure setup and data
        distribution as \Cref{fig:ve_gmm_denoising}. Note that compared
        to \Cref{fig:ve_gmm_denoising}, the noise distribution is much
    more concentrated around the origin.}
    \label{fig:vp_gmm_denoising}
\end{figure}

\paragraph{Step 2: Different noise models.} The second step is to
consider slightly different models compared to
\eqref{eq:additive_gaussian_noise_model}. The basic motivation for
this is as follows. In practice, the noise distribution
\(\dNorm(\vzero, t_{L}^{2}\vI)\) becomes an increasingly poor
estimate of the true covariance in high dimensions, i.e.,
\eqref{eq:ve_converge_to_large_gaussian} becomes an increasingly
poor approximation, especially with anisotropic high-dimensional
data. The increased distance between \(\dNorm(\vzero, t_{L}^{2}\vI)\)
and the true distribution of \(\vx_{t_{L}}\) may cause the denoiser
to perform worse in such circumstances. Theoretically,
\(\vx_{t_{L}}\) never converges to any distribution as \(t_{L}\)
increases, so this setup is difficult to analyze end-to-end. In this
case, our remedy is to \textit{simultaneously add noise and shrink
    the contribution of \(\vx\), such that \(\vx_{T}\) converges as \(T
\to \infty\)}. The rate of added noise is denoted \(\sigma \colon [0,
T] \to \R_{\geq 0}\), and the rate of shrinkage is denoted \(\alpha
\colon [0, T] \to \R_{\geq 0}\), such that \(\sigma\) is
\textit{increasing} and \(\alpha\) is (not strictly) \textit{decreasing}, and
\begin{equation}\label{eq:gen_additive_gaussian_noise_model}
    \vx_{t} \doteq \alpha_{t}\vx + \sigma_{t}\vg, \qquad \forall t \in [0, T].
\end{equation}
The previous setup has \(\alpha_{t} = 1\) and \(\sigma_{t} = t\), and
this is called the \textit{variance-exploding (VE) process}. A
popular choice that decreases the contribution of \(\vx\), as we
described originally, has \(T = 1\) (so that \(t \in [0, 1]\)),
\(\alpha_{t} = \sqrt{1 - t^{2}}\) and \(\sigma_{t} = t\); this is the
\textit{variance-preserving (VP) process}. Note that under the VP
process, \(\vx_{1} \sim \dNorm(\vzero, \vI)\) exactly, so we can just
sample from this standard distribution and iteratively denoise. As a
result, the VP process is much easier to analyze theoretically and
more stable empirically.\footnote{Why use the whole \(\alpha,
    \sigma\) setup? As we will see in
    \Cref{exercise:implement_denoising_processes}, it encapsulates and
    unifies many proposed processes, including the recently popular
    so-called \textit{flow matching} process. Despite this, the VE and VP
    processes are still the most popular empirically and theoretically
(so far), and so we will consider them in this Section.}

With this more general setup, Tweedie's formula \eqref{eq:tweedie} becomes
\begin{equation}\label{eq:gen_tweedie}
    \Ex[\vx \mid \vx_{t}] = \frac{1}{\alpha_{t}}\bp{\vx_{t} +
    \sigma_{t}^{2}\nabla \log p_{t}(\vx_{t})}.
\end{equation}
The denoising iteration \eqref{eq:denoising_iteration} becomes
\begin{equation}\label{eq:gen_denoising_iteration}
    \hat{\vx}_{t_{\ell - 1}} = \frac{\sigma_{t_{\ell -
    1}}}{\sigma_{t_{\ell}}}\hat{\vx}_{t_{\ell}} + \bp{\alpha_{t_{\ell
        - 1}} - \frac{\sigma_{t_{\ell -
    1}}}{\sigma_{t_{\ell}}}\alpha_{t_{\ell}}}\bar{\vx}^{\star}(t_{\ell},
    \hat{\vx}_{t_{\ell}}).
\end{equation}

\begin{example}[Gaussian mixture model]
    For a Gaussian mixture model (see \Cref{ex:gaussian_mixture}),
    its denoiser \eqref{eq:gmm_bayes_optimal_denoiser} becomes
    \begin{align}\label{eq:gen_gmm_bayes_optimal_denoiser}
        \bar{\vx}^{\star}(t, \vx_{t})
        &= \sum_{k = 1}^{K}\Bigg\{\frac{\pi_{k}\phi(\vx_{t}
                    ; \alpha_{t}\vmu_{k}, \alpha_{t}^{2}\vSigma_{k}
            + \sigma_{t}^{2}\vI)}{\sum_{i = 1}^{K}\pi_{i}\phi(\vx_{t} ;
                    \alpha_{t}\vmu_{i}, \alpha_{t}^{2}\vSigma_{i} +
            \sigma_{t}^{2}\vI)} \\
            &\qquad \qquad \qquad \cdot\bp{\vmu_{k} +
                \alpha_{t}\vSigma_{k}(\alpha_{t}^{2}\vSigma_{k} +
        \sigma_{t}^{2}\vI)^{-1}(\vx_{t} - \alpha_{t}\vmu_{k})}\Bigg\}. \nonumber
    \end{align}
    \Cref{fig:vp_gmm_denoising} demonstrates iterations of the
    sampling procedure.
\end{example}

Note that the denoising iteration \eqref{eq:gen_denoising_iteration}
gives a sampling algorithm called the DDIM (``Denoising Diffusion
Implicit Model'') sampler \cite{song2020denoising}, and is one of the
most popular sampling algorithms used in diffusion models today. We
summarize it in \Cref{alg:iterative_denoising}.

\begin{algorithm}
    \caption{Sampling using a denoiser.}
    \label{alg:iterative_denoising}
    \begin{algorithmic}[1]
        \Require{An ordered list of timesteps \(0 \leq t_{0} < \cdots
        < t_{L} \leq T\) to use for sampling.}
        \Require{A denoiser \(\bar{\vx} \colon \{t_{\ell}\}_{\ell =
        1}^{L} \times \R^{D} \to \R^{D}\).}
        \Require{Scale and noise level functions \(\alpha, \sigma
        \colon \{t_{\ell}\}_{\ell = 0}^{L} \to \R_{\geq 0}\).}
        \Ensure{A sample \(\hat{\vx}\), approximately from the
        distribution of \(\vx\).}
        \Function{DDIMSampler}{$\bar{\vx}, (t_{\ell})_{\ell = 0}^{L}$}
        \State{Initialize \(\hat{\vx}_{t_{L}} \sim\) approximate
        distribution of \(\vx_{t_{L}}\)}
        \Statex{\quad\ \texttt{\# VP \(\implies
        \dNorm(\vzero, \vI)\), VE \(\implies \dNorm(\vzero, t_{L}^{2}\vI)\).}}
        \Statex{}
        \For{\(\ell = L, L - 1, \dots, 1\)}
        \State{Compute
            \begin{equation*}
                \hat{\vx}_{t_{\ell - 1}} \doteq \frac{\sigma_{t_{\ell
                - 1}}}{\sigma_{t_{\ell}}}\hat{\vx}_{t_{\ell}} +
                \bp{\alpha_{t_{\ell - 1}} - \frac{\sigma_{t_{\ell -
                1}}}{\sigma_{t_{\ell}}}\alpha_{t_{\ell}}}\bar{\vx}(t_{\ell},
                \hat{\vx}_{t_{\ell}})
            \end{equation*}
        }
        \EndFor
        \State{\Return{\(\hat{\vx}_{t_{0}}\)}}
        \EndFunction
    \end{algorithmic}
\end{algorithm}

\paragraph{Step 3: Optimizing training pipelines.} If we use the
procedure dictated by \Cref{sub:intro_diffusion_denoising} to learn a
separate denoiser \(\bar{\vx}(t, \cdot)\) for each time \(t\) to be
used in the sampling algorithm, \textit{we would have to learn \(L\)
separate denoisers!} This is highly inefficient---the usual case is
that we have to train \(L\) separate neural networks, taking up \(L\)
times the training time and storage memory, and then be locked into
using these timesteps for sampling forever. Instead, we can
\textit{train a single neural network} to denoise across all times
\(t\), taking as input the continuous variables \(\vx_{t}\) and \(t\)
(instead of just \(\vx_{t}\) before). Mechanically, our training loss
averages over \(t\), i.e., solves the following problem:
\begin{equation}\label{eq:training loss}
    \min_{\theta}\Ex\norm{\bar{\vx}_{\theta}(\tau,
    \vx_{\tau}) - \vx}_{2}^{2},
\end{equation}
where \(\tau \sim \dUnif([0, T])\), i.e., is drawn uniformly from \([0, T]\).
Similar to Step 1, where we used more timesteps closer to \(t = 0\)
to ensure a better sampling process, we may want to ensure that the
denoiser is higher quality closer to \(t = 0\), and thereby
\textit{weight the loss} so that \(t\) near \(0\) has higher
weight.\footnote{In practice, the distribution of sampled \(t\) may
    also be non-uniform, in order to ensure (in conjunction with the
    weighting scheme) that the Monte Carlo estimate for the loss has
small variance, greatly improving training stability.}
Letting \(w_{t}\) be the weight at time \(t\), the weighted loss would look like
\begin{equation}\label{eq:parametric_denoising_loss}
    \min_{\theta}\Ex[w_{\tau}\norm{\bar{\vx}_{\theta}(\tau,
    \vx_{\tau}) - \vx}_{2}^{2}].
\end{equation}
One reasonable choice of weight in practice is \(w_{t} =
\alpha_{t}^{2}/\sigma_{t}^{2}\). The precise reason will be covered in the
next paragraph, but generally it serves to up-weight the losses
corresponding to \(t\) near \(0\) while still remaining reasonably
numerically stable. Also, of course, we cannot compute the
expectation in practice, so we use the most straightforward
Monte Carlo average to estimate it. The series of changes made here
have several conceptual and computational benefits: we do not need to
train multiple denoisers, we can train on one set of timesteps and
sample using a subset (or others entirely), etc. The full pipeline is
discussed in \Cref{alg:learning_denoiser}. Certainly, this pipeline
works in practice; we discuss some results in
\Cref{sec:auto-encode-img-gen,sec:cond_gen,sec:object-generation,sec:shape-and-pose,sec:egoallo}.

\begin{algorithm}
    \caption{Learning a denoiser from data.}
    \begin{algorithmic}[1]
        \Require{Dataset \(\cD \subseteq \R^{D}\).}
        \Require{An ordered list of timesteps \(0 \leq t_{0} < \cdots
        < t_{L} \leq T\) to use for sampling.}
        \Require{A weighting function \(w \colon \{t_{\ell}\}_{\ell =
        1}^{L} \to \R_{\geq 0}\).}
        \Require{Scale and noise level functions \(\alpha, \sigma
        \colon \{t_{\ell}\}_{\ell = 0}^{L} \to \R_{\geq 0}\).}
        \Require{A parameter space \(\Theta\) and a denoiser
        architecture \(\bar{\vx}_{\theta}\).}
        \Require{An optimization algorithm for the parameters.}
        \Require{The number of optimization iterations \(M\).}
        \Require{The number of Monte Carlo draws \(N\) per training step.}
        \Ensure{A trained denoiser \(\bar{\vx}_{\theta^{\star}}\).}

        \Function{TrainDenoiser}{$\cD, \Theta$}
        \State{Initialize \(\theta^{(1)} \in \Theta\)}
        \For{\(m \in [M]\)}
        \For{\(i \in [N]\)}
        \Statex{\qquad\qquad\ \texttt{\# Draw a sample from
        the dataset.}}
        \State{\(\vx^{(m, i)} \sim \cD\)}
        \Statex{}

        \Statex{\qquad\qquad\ \texttt{\# Sample a timestep.}}
        \State{\(t^{(m, i)} \simiid \dUnif(\{t_{\ell}\}_{\ell =
        1}^{L})\)}
        \Statex{}

        \Statex{\qquad\qquad\ \texttt{\# Compute weights and scales.}}
        \State{\((\alpha^{(m, i)}, \sigma^{(m, i)}, w^{(m, i)}) =
        (\alpha_{t^{(m, i)}}, \sigma_{t^{(m, i)}}, w_{t^{(m, i)}})\)}
        \Statex{}

        \Statex{\qquad\qquad\ \texttt{\# Sample a noise vector.}}
        \State{\(\vg^{(m, i)} \simiid \dNorm(\vzero, \vI)\)}
        \Statex{}

        \Statex{\qquad\qquad\ \texttt{\# Compute the noised sample.}}
        \State{\(\vx_{t}^{(m, i)} \doteq \alpha^{(m, i)}\vx^{(m, i)}
        + \sigma^{(m, i)}\vg^{(m, i)}\)}
        \EndFor

        \Statex{\qquad\ \texttt{\# Compute the training loss.}}
        \State{\(\hat{\cL}^{(m)} \doteq \displaystyle
                \frac{1}{N}\sum_{i = 1}^{N}w^{(m, i)}\norm{\vx^{(m, i)} -
                    \bar{\vx}_{\theta^{(m)}}(t^{(m, i)}, \vx_{t}^{(m,
        i)})}_{2}^{2}\)}
        \Statex{}
        \Statex{\qquad\ \texttt{\# Update the model parameters.}}
        \State{\(\theta^{(m + 1)} \doteq
                \texttt{OptimUpdate}^{(m)}(\theta^{(m)},
        \nabla_{\theta^{(m)}}\hat{\cL}^{(m)})\)}
        \EndFor
        \Statex{\quad\ \texttt{\# Return the trained denoiser.}}
        \State{\Return{\(\bar{\vx}_{\theta^{(M + 1)}}\)}}
        \EndFunction
    \end{algorithmic}
    \label{alg:learning_denoiser}
\end{algorithm}

\paragraph{Step 4: Changing the estimation target.}
It is common to instead reorient the whole denoising pipeline
around \textit{noise predictors}, i.e., estimates of \(\Ex[\vg \mid
\vx_{t}]\). In practice, noise predictors are slightly easier to
train because their output is (almost) always of comparable size to a
Gaussian random variable, so training is more numerically stable, but
they may require larger models since they predict a
full-dimensional target instead of potentially low-dimensional data
\citep{li2025back}. By
\eqref{eq:gen_additive_gaussian_noise_model} we have
\begin{equation}
    \vx_{t} = \alpha_{t}\Ex[\vx \mid \vx_{t}] + \sigma_{t}\Ex[\vg
    \mid \vx_{t}] \implies \Ex[\vg \mid \vx_{t}] =
    \frac{1}{\sigma_{t}}\bp{\vx_{t} - \alpha_{t}\Ex[\vx \mid \vx_{t}]},
\end{equation}
so any predictor for \(\vx\) can be turned into a predictor
for \(\vg\) using the above relation, i.e.,
\begin{equation}
    \bar{\vg}(t, \vx_{t}) = \frac{1}{\sigma_{t}}\vx_{t} -
    \frac{\alpha_{t}}{\sigma_{t}}\bar{\vx}(t, \vx_{t}),
\end{equation}
and vice-versa. Thus a good network for estimating \(\bar{\vg}\) is
the same as a good network for estimating \(\bar{\vx}\) \textit{plus
a residual connection} (as seen in, e.g., transformers). Their losses
are also the same as the denoiser, up to the factor of
\(\alpha_{t}/\sigma_{t}\), i.e.,
\begin{equation}
    \Ex[w_{\tau}\norm{\vg - \bar{\vg}(\tau,
    \vx_{\tau})}_{2}^{2}] =
    \Ex\rs{w_{\tau}\frac{\alpha_{\tau}^{2}}{\sigma_{\tau}^{2}}\norm{\vx
    - \bar{\vx}(\tau, \vx_{\tau})}_{2}^{2}},
\end{equation}
recalling that \(\tau \sim \dUnif([0, T])\). Other targets have
been proposed for different tasks, e.g., \(\Ex[\vv_{t} \mid
\vx_{t}]\) where \(\vv_{t} = \odv*{\vx_{t}}{t}\) (called
\(v\)-prediction or velocity prediction), etc. This turns out to be
equivalent to denoising and noise prediction, but also fundamentally
useful for generalizations of diffusion such as flow matching, which
we will discuss in \Cref{rem:flow_matching}.

\subsubsection{Sampling Error Rate of the Denoising Process}
We have made many changes to our original platonic noising/denoising
process. To assure ourselves that the new process still works in
practice, we can compute numerical examples (such as
\Cref{fig:vp_gmm_denoising}). To assure ourselves that it is
theoretically sound, we can prove a \textit{bound on the error rate}
for the sampling algorithm, which shows that the error rate is small.
We now furnish such a rate from the literature, which shows that the
output distribution of the sampler converges in the so-called
\textit{total variation (TV) distance} to the true distribution. The
TV distance between two random variables \(\vx\) and
\(\vy\) is defined as:\footnote{The supremum here is only over
    \textit{measurable} sets \(A\), not absolutely all sets. As with all
    other measure-theoretic caveats in this book, feel free to
ignore it for practical purposes.}
\begin{equation}
    \TV(\vx, \vy) \doteq \sup_{A \subseteq \R^{d}}\abs*{\Pr[\vx \in
    A] - \Pr[\vy \in A]}.
\end{equation}
If \(\vx\) and \(\vy\) are very close (uniformly), then the supremum
will be small. So the TV distance measures the closeness of random
variables. (It is indeed a metric, as the name suggests; the proof is
an exercise.)

\begin{theorem}[\citep{li2024d} Theorem 1,
    Simplified]\label{thm:diffusion_sampler_convergence}
    Suppose that \(\Ex\norm{\vx}_{2} < \infty\).  If \(\vx\) is
    denoised according to the VP process with an exponential
    discretization\footnote{The precise definition is rather lengthy
        in our notation and only defined up to various absolute
        constants, so we omit it here for brevity. It is in the
    original paper \citep{li2024d}.} as in
    \eqref{eq:denoising_exponential_discretization}, the output
    \(\hat{\vx}\) of \Cref{alg:iterative_denoising} satisfies the
    total variation bound
    \begin{align}\label{eq:diffusion_sampling_error}
        \TV(\vx, \hat{\vx})
        &= \tilde{\cO}\Bigg(\underbrace{\frac{D}{L}}_{\text{discretization
            error}} \\
            & \qquad \qquad + \underbrace{\sqrt{\frac{1}{L}\sum_{\ell =
                    1}^{L}\frac{\alpha_{t_{\ell}}^{2}}{\sigma_{t_{\ell}}^{4}}\Ex\norm{\bar{\vx}^{\star}(t_{\ell},
                        \vx_{t_{\ell}}) - \bar{\vx}(t_{\ell},
            \vx_{t_{\ell}})}_{2}^{2}}}_{\text{average excess error of
        the denoiser}}\Bigg)\nonumber
    \end{align}
    where \(\bar{\vx}^{\star}\) is the Bayes optimal denoiser for
    \(\vx\), and \(\tilde{\cO}\) is a version of the big-\(\cO\)
    notation that ignores logarithmic factors in \(L\).
\end{theorem}
The very high-level proof technique is, as discussed earlier, to
bound the error at each step, distinguish the error sources (between
discretization and denoiser error), and carefully ensure that the
errors do not accumulate too much (or even cancel out).

Note that if \(L \to \infty\) and we correctly learn the Bayes-optimal
denoiser \(\bar{\vx} = \bar{\vx}^{\star}\) (so the excess error is
\(0\)), then the sampling process in
\Cref{alg:iterative_denoising} yields a \textit{perfect (in
distribution) inverse} of the noising process, since the error rate in
\Cref{thm:diffusion_sampler_convergence} goes to
\(0\),\footnote{Similar results hold for VE processes; the theoretical
    toolkits used for the proofs differ because of the structure of the
underlying stochastic processes.} as argued heuristically above.

\begin{remark}
    What if the data is low-dimensional, say supported on a
    \textit{low-rank subspace} of the high-dimensional space
    \(\R^{D}\)? If the data distribution is compactly supported---say
    if the data is normalized to the unit hypercube, which is often
    ensured as a preprocessing step for real data such as
    images---it is possible to do better. Namely, the authors of
    \cite{li2024d} also define a measure of \textit{approximate
    intrinsic dimension} \(\kappa\) using the asymptotics of the so-called
    covering number, which is extremely similar in intuition (if not
    in implementation) to the rate-distortion function presented in
    the next Chapter. Then they show that using a particular small
    modification of the DDIM sampler in
    \Cref{alg:iterative_denoising} (i.e., slightly perturbing the
    update coefficients), the discretization error becomes
    \(\tilde{\cO}(\kappa/L)\) instead of \(\tilde{\cO}(D/L)\) as in
    \Cref{thm:diffusion_sampler_convergence}. Therefore, using this
    modified algorithm, \(L\) does not have to be too large even as
    \(D\) reaches the thousands or millions, since real data have
    low-dimensional structure. However, in practice we use the DDIM
    sampler instead, so \(L\) should have a mild dependence on \(D\)
    to achieve consistent error rates. The exact choice of \(L\)
    trades off between the computational complexity (e.g., runtime or
    memory consumption) of sampling and the statistical complexity of
    learning a denoiser for low-dimensional structures. The value of
    \(L\) is often different at training time (where a larger \(L\)
        allows better coverage of the interval \([0, T]\), which helps
    the network learn a relationship that generalizes over \(t\))
    and sampling time (where \(L\) being smaller means more efficient
    sampling). One can even pick the timesteps adaptively at sampling
    time to optimize this tradeoff \cite{bao2022analytic}.
\end{remark}

\begin{remark}
    Various other works define the reverse process as moving backward
    in the time index \(t\) using an explicit difference equation, or
    differential equation in the limit \(L \to \infty\), or forward
    in time using the transformation \(\vy_{t} = \vx_{T - t}\), such
    that if \(t\) increases then \(\vy_{t}\) becomes closer to
    \(\vx_{0}\). Each paper uses its own notation, and it can
    sometimes be confusing to understand the relationships between
    different papers' results. In this work we strive to keep
    consistency: we move forward in time to noise, and backward in
    time to denoise. If you are reading another work that is not
    clear on the time index, or trying to implement an algorithm
    that is similarly unclear, there is one way to disambiguate
    every time: the sampling process should always have a
    \textit{positive} coefficient on the denoiser term (or
    equivalently the score function) when taking a step.
\end{remark}

\begin{remark}\label{rem:flow_matching}
    The diffusion/denoising processes and associated models provide a
    way to transport the data distribution to and from a high-entropy
    (Gaussian or Gaussian-like) distribution. On the other hand, it
    may be useful to transport the data distribution to and from an
    arbitrary target distribution, which could be Gaussian, a
    Gaussian mixture, or even another type of data distribution. If
    we have this desideratum, we can use the \textit{flow matching}
    framework to achieve this. Suppose the target random variable is
    \(\vy\), and our flow is
    \begin{equation}
        \vx_{t} = \alpha_{t}\vx + \sigma_{t}\vy, \qquad \forall t \in [0, 1].
    \end{equation}
    such that \(\alpha_{0} = \sigma_{1} = 1\) and \(\alpha_{1} =
    \sigma_{0} = 0\), so that \(\vx_{0} = \vx\) and \(\vx_{1} =
    \vy\). Nothing about our strategy changes: we
    can train a denoiser \(\bar{\vx}_{\theta}\) to denoise
    \(\vx_{t}\) back to \(\vx\) (or equivalently a ``noise''
    predictor to predict \(\vy\)), and then apply our iterative
    denoising algorithms to start from a realization of \(\vy\) to
    obtain \(\vx\). The only thing missing is Tweedie's formula which
    connects the score \(\nabla \log p_{t}\) of the denoiser to the
    denoiser; such a formula does not exist in general.

    In practice, inspired by the concept of learning a so-called
    \textit{transport map} between data pairs \(\vx\) and \(\vy\),
    practitioners often wish to train a \textit{velocity predictor}
    \(\bar{\vv}_{\theta}\) to predict \(\vv_{t} = \odv*{\vx_{t}}{t} =
    \alpha_{t}^{\prime}\vx + \sigma_{t}^{\prime}\vy\) instead of a
    denoiser \(\bar{\vx}_{\theta}\), or rather train
    \(\bar{\vv}_{\theta}\) to predict \(\Ex[\vv_{t} \mid \vx_{t}]\).
    In this case, it is simplest to take the choice of coefficients
    \(\alpha_{t} = 1 - t\) and
    \(\sigma_{t} = t\) (for \(t \in [0, 1]\)), so that the prediction
    target is \(\vv_{t} = \vy - \vx\) for each time \(t\). Given the
    estimator \(\bar{\vv}_{\theta}\), the standard sampling algorithm
    follows the field backwards in time, i.e., starting from a
    realization of \(\vy\), we can iteratively update the iterate as follows:
    \begin{equation}
        \hat{\vx}_{t_{\ell - 1}} = \hat{\vx}_{t_{\ell}} -
        \bar{\vv}_{\theta}(t_{\ell}, \hat{\vx}_{t_{\ell}}) \Delta t,
    \end{equation}
    where \(\Delta t = t_{\ell} - t_{\ell - 1}\) is the timestep.
    Note that velocity prediction is algebraically equivalent to
    denoising (see \Cref{exercise:velocity_prediction}). See
    \Cref{sec:auto-encode-img-gen,sec:shape-and-pose} for practical
    applications of flow matching.
\end{remark}

\begin{remark}
    In practice, diffusion usually takes many steps and many function
    evaluations of the denoiser. It is therefore natural to ask:
    ``can we speed up the sampling process by using a smaller number
    of steps and function evaluations?'' The answer is yes, and such
    models are often called \textit{consistency models} or, more
    descriptively, \textit{few-step models}. The core idea is to use a
    denoiser that takes in two times \(\bar{\vx}_{\theta}(s, t,
    \vx_{t})\) and estimates \(\Ex[\vx_{s} \mid \vx_{t}]\), which in
    some sense is
    an easier task, or alternatively to use a velocity predictor
    \(\bar{\vv}_{\theta}(s, t, \vx_{t})\) and estimate the average
    velocity \(\Ex[(\vx_{t} - \vx_{s})/(t - s) \mid \vx_{t}]\). Such
    denoisers have additional consistency
    conditions such as \(\bar{\vx}_{\theta}(u, t,
    \bar{\vx}_{\theta}(s, u, \vxi)) = \bar{\vx}_{\theta}(s, t, \vxi)\),
    which motivate auxiliary training losses. Once such a predictor
    is learned, the
    sampling process becomes very simple and efficient, as it only
    requires a few steps of the form
    \begin{equation}
        \hat{\vx}_{s} = \bar{\vx}_{\theta}(s, t, \hat{\vx}_{t}) =
        \hat{\vx}_{t} - \bar{\vv}_{\theta}(s, t, \hat{\vx}_{t})\cdot(t - s).
    \end{equation}
    Relevant formulations include flow maps \cite{boffi2025build} and
    mean flows \cite{geng2025mean}.
\end{remark}

\section{Memorization, Generalization, and Coding Rates}
\label{sec:diffusion_memorization}

The calculation presented at the end of
\Cref{sub:intro_diffusion_denoising} seems to suggest (loosely
speaking) that in practice, using a transformer-like network is a
good choice for learning or approximating a denoiser. This is
reasonable, but what is the problem with using any old neural network
(such as a multi-layer perceptron (MLP)) and simply scaling it
up to infinity? To answer this question, let us consider a constraint
we have ignored or glossed over so far in this Chapter: we only
have finite data with which to learn the denoiser. Indeed, let us
consider the training problem \eqref{eq:parametric_denoising_loss},
except this time we assume that our data \(\vx\) are drawn from the
empirical distribution \(p^{\vX}\) on \(N\) training samples
\(\vX = [\vx^{(1)}, \dots, \vx^{(N)}]\). Because \(p^{\vX}\) is supported
on a \(0\)-dimensional subset of the ambient space \(\R^{D}\), we
have that \(h(p^{\vX}) = -\infty\), so in some sense \(p^{\vX}\) is
the simplest (i.e., lowest-entropy) distribution that could generate
our observations \(\vX\). In short, this assumption arguably \textit{has the
least inductive bias} about the data distribution! Many
popular commentaries on deep learning and generative models would
have you believe that, due to this property, the best way to solve
the generative modeling problem is to scale a generic denoiser
as large as possible and reap the corresponding
benefits. However, this advice is (provably!) wrong in this setting.
To understand why, let us note that the empirical distribution
\(p^{\vX}\) is actually a mixture of Gaussians with zero covariances:
\begin{equation}
    p^{[\vx^{(1)}, \dots, \vx^{(N)}]} = \frac{1}{N}\sum_{i =
    1}^{N}\dNorm(\vx^{(i)}, \vzero).
\end{equation}
As such, we know that the (globally) optimal solution to the
denoising problem \eqref{eq:parametric_denoising_loss} with \(\vx\)
drawn from the empirical distribution, written out as follows:
\begin{equation}\label{eq:parametric_denoising_loss_finite_sample}
    \min_{\theta \in \Theta}\frac{1}{N}\sum_{i =
    1}^{N}\Ex[w_{\tau}\norm{\vx^{(i)} - \bar{\vx}_{\theta}(\tau,
    \vx_{\tau}^{(i)})}_{2}^{2}]
\end{equation}
where \(\vx_{t}^{(i)}\) is the noised version of the training sample
\(\vx^{(i)}\). The optimizer of this loss over all
(square-integrable) denoisers \(\bar{\vx}\) is the associated optimal
Gaussian denoiser from \Cref{example:denoising_gaussian_mixture}, i.e.,
\begin{equation}\label{eq:memorizing_denoiser}
    \bar{\vx}^{\star}(t, \vx_{t}) = \sum_{i =
    1}^{N}\frac{\exp({-\|\vx_{t} -
    \alpha_{t}\vx^{(i)}\|_{2}^{2}/(2\sigma_{t}^{2})})}{\sum_{j =
        1}^{N}\exp({-\|\vx_{t} -
    \alpha_{t}\vx^{(j)}\|_{2}^{2}/(2\sigma_{t}^{2})})}\vx^{(i)}.
\end{equation}
This is a convex combination of the data \(\vx^{(i)}\), and the
coefficients get ``sharper'' (i.e., closer to \(0\) or \(1\)) as \(t
\to 0\). In particular, for small \(t\), the denoiser is almost
exactly equal to the closest sample in the training set, and the
update moves the iterate closer to the closest point, such that at
the end of sampling, the iterate \(\hat{\vx}\) is (nearly) equal to a
training sample  \(\vx^{(i)}\). Let us re-emphasize this point:
\textit{if we use the denoiser \eqref{eq:memorizing_denoiser} and a
good sampler, our samples will re-generate the training data}. This
is confirmed by our theoretical guarantee on sampling
\Cref{thm:diffusion_sampler_convergence}, which, when instantiated in
this context, says that if we take the number of sampling steps to
\(\infty\) and use the prescribed timestep schedule, the diffusion
model associated with the denoiser \eqref{eq:memorizing_denoiser}
will always sample from the empirical distribution and thus always
reproduce a training sample. This phenomenon is a particularly strict
instantiation of \textit{memorization} in generative
modeling.\footnote{Memorization is a much broader phenomenon whose
    exact characterization depends on the domain the model is trained on
    (e.g., text, images, videos, audio, etc.). Later in this Section
we will discuss more heuristic measures of memorization.}

Since in practice we only have a fixed finite dataset on which we can
train our denoiser, the learning problem we aim to solve in this case
is the finite-sample denoising loss
\eqref{eq:parametric_denoising_loss_finite_sample} (compared to the
    infinite-sample denoising loss \eqref{eq:parametric_denoising_loss},
for example). As such, \textit{if we have an arbitrarily large model
    and optimize it as well as possible, it will learn to memorize and
re-generate the training set}. In other words, this very large-scale
model is not much more interesting than a simple database of the training data.

\begin{figure}
    \centering
    \includegraphics[width=\textwidth]{\toplevelprefix/chapters/chapter3/figs/euclidean_memorization_samples.png}
    \caption{\small\textbf{Generated samples (left) from a diffusion
            model trained on \(6{,}000\) samples and (right) their closest
        points in the training dataset (in the Euclidean norm).} We
        observe that the generated samples are essentially pixel-perfect
        equivalents to the most similar training data. Both sets of images
    are from \cite{yoon2023diffusion}.}
    \label{fig:diffusion_memorization_euclidean_samples}
\end{figure}

\begin{figure}
    \centering
    \includegraphics[width=\textwidth]{\toplevelprefix/chapters/chapter3/figs/euclidean_memorization_ratio.png}
    \caption{\small\textbf{Memorization versus dataset size for a
            fixed model, on (left) 4,000 training epochs and (right)
        40,000 training epochs with smaller datasets}. We observe
        from the left plot that large enough datasets will never be
        memorized (e.g., 10,000 to 50,000 samples), whereas small
        datasets (1,000 to 5,000 samples) will be memorized after
        enough training. This means that as the dataset becomes
        larger relative to the model, the memorization ratio at the
        end of training goes from near-1 to near-0 relatively
        rapidly, indicating the possible existence of a phase
        transition from memorization to generalization as the dataset
        size increases. Figure taken from
    \cite{gu2023memorization}.}
\end{figure}

The above argument holds for arbitrarily powerful models that can
fit the (true) memorizing denoiser. It is therefore relevant to ask
whether it actually bears out in practice: do very large models
(relative to the dataset) learn to memorize the training data? It
turns out that this is empirically true. One of the first studies on
this was \cite{yoon2023diffusion}. To understand their results, we define the
memorization ratio as the probability of memorization:
\begin{equation}
    \mathrm{MemRatio}(\bar{\vx}) = \Pr_{\hat{\vx} \sim
    \mathrm{DM}(\bar{\vx})}[\hat{\vx}\ \text{is memorized}],
\end{equation}
where \(\mathrm{DM}(\bar{\vx})\) is the distribution induced by the
diffusion model that uses denoiser \(\bar{\vx}\), and the sample
\(\hat{\vx}\) is memorized using the following criterion:
\begin{equation}\label{eq:diffusion_euclidean_memorization_criterion}
    \hat{\vx}\ \text{is memorized}\ \iff \frac{\norm{\hat{\vx} -
    \vx_{(1)}(\hat{\vx})}_{2}}{\norm{\hat{\vx} -
    \vx_{(2)}(\hat{\vx})}_{2}} < \frac{1}{3},
\end{equation}
where \(\vx_{(i)}(\hat{\vx})\) is the \(i\th\) closest point to
\(\hat{\vx}\) in the training dataset (in the Euclidean norm), and
\(1/3\) is a constant that is empirically shown to agree well with
human perception of memorization. Notice that this definition exactly
corresponds to our earlier intuitive definition of memorization as
sampled points each being much closer to a particular training sample than all
other training samples. They showed that this kind of memorization
occurs in practice when the model is much larger (in terms of the
number of parameters) relative to the data;
\Cref{fig:diffusion_memorization_euclidean_samples} demonstrates that
in this case, generated samples are essentially pixel-perfect copies
of training samples. Slightly later, \cite{gu2023memorization}
performed a careful study to characterize how various factors such as
architectures and training time impact memorization, showing among
other things that as the dataset size increases relative to a fixed
model, the memorization ratio rapidly changes from near-1 to near-0,
i.e., the model quickly stops memorizing. This
indicates the possibility of a phase transition between memorization
and generalization as the model becomes smaller relative to the data.
\cite{buchanan2025edge} investigates the existence and theoretical
mechanism of this phase transition in synthetic data drawn from a
Gaussian mixture model, demonstrating that it exists and providing a partial
characterization of the model size at which models will start
or stop memorizing. Other theoretical work such as
\cite{niedoba2024towards,kamb2024analytic,george2025denoising,zhang2025generalization}
also attempts to provide a theoretical characterization of the phase
transition in terms of model and dataset size in different settings.

\begin{figure}
    \centering
    \includegraphics[width=0.475\textwidth]{\toplevelprefix/chapters/chapter3/figs/generalizability.png}
    \hfill
    \includegraphics[width=0.45\textwidth]{\toplevelprefix/chapters/chapter3/figs/generalizability_loss.png}
    \caption{\small\textbf{The generalization score (left) and
            training loss (right) for a
        fixed model and increasing data}, from
        \cite{pmlr-v235-zhang24cn}. The generalization score is \(1 - \)
        the fraction of samples that are memorized, and the different
        lines correspond to differently-sized denoiser models
        (UNet-64 is smaller than UNet-128, which in turn is smaller
        than UNet-256). The figures show that, as the dataset size
        becomes larger relative to the model size, the loss increases
        significantly and memorization stops occurring, showing that our
    theory and intuition apply to real models.}
    \label{fig:diffusion_generalizability_loss}
\end{figure}

In order to empirically prove the existence of a phase transition
between memorization and generalization in real data,
\cite{pmlr-v235-zhang24cn} rethought the definition of memorization
using the following reasoning. In the very high-dimensional space of
natural images, Euclidean distances between points become much more
strict \cite{Wright-Ma-2022}, in the sense that if two images are not
pixel-perfect matches to each other, then their Euclidean distances
will concentrate sharply around a fixed value (scaling with the
images' dimension). This does not catch the case where certain
\textit{concepts} (such as copyrighted cartoon characters) are
generated in novel positions, making it very difficult to form
effective and practical memorization criteria around the Euclidean
distance. To fix this, \cite{pmlr-v235-zhang24cn} uses a particular
function \(f_{\mathrm{SSCD}}\) (a pre-trained neural
network called a Self-Supervised Copy Detector \cite{pizzi2022self})
to map images into a Euclidean space, then takes the normalized inner
product to obtain their (cosine) similarity, such that images that
are similar due to this cosine similarity are claimed to be
conceptually similar in the sense needed for memorization. Then,
\cite{pmlr-v235-zhang24cn} says that a sample \(\hat{\vx}\) is memorized if its
similarity with any training point is \(> 3/5\), i.e.,
\begin{equation}\label{eq:diffusion_sscd_memorization_criterion}
    \hat{\vx}\ \text{is memorized} \iff \max_{i \in [N]}
    \frac{f_{\mathrm{SSCD}}(\hat{\vx})^{\top}f_{\mathrm{SSCD}}(\vx^{(i)})}{\norm{f_{\mathrm{SSCD}}(\hat{\vx})}_{2}\norm{f_{\mathrm{SSCD}}(\vx^{(i)})}_{2}}
    > \frac{3}{5}
\end{equation}

\Cref{fig:diffusion_generalizability_loss} shows
the outcome of taking many (around 10,000) samples from diffusion
models trained on differently-sized training datasets in terms of
estimating memorization and computing the training loss.
The outcome is in line with the above discussion; when the model is
too large relative to the training data and the training loss is too small,
memorization occurs and samples reproduce the training data. Thus, our
intuition applies in practice: \textit{for generative diffusion models, scaling
the model is \underline{not} all you need!}

So now the question is: what kind of denoiser should we use? The key
is to choose a network architecture that can approximate the true
denoiser (say, the one corresponding to a low-rank distribution as in
\eqref{eq:gmm_lowrank_denoiser}) without approximating the memorizing
denoiser \eqref{eq:memorizing_denoiser}. Some work has explored this
fine line and why modern diffusion models, which use Transformer- and
convolution-based architectures, can memorize and generalize in
different regimes
\citep{kamb2024analytic,niedoba2024towards,zhang2025generalization}.

\begin{figure}
    \centering
    \includegraphics[width=0.4\textwidth]{\toplevelprefix/chapters/chapter3/figs/diffusion_perturbation.png}
    \caption{\small\textbf{Visualizing the action of a well-trained
        denoiser} as a perturbation of the empirical denoiser at
        \textit{small times}, when the memorizing denoiser acts like a
    projection onto the nearest point in the training set.}
    \label{fig:diffusion_perturbation}
\end{figure}

Let us at last return to the initial premise of this Chapter, described
in \Cref{sub:min_entropy}: finding the minimal-entropy distribution
that models the data is the only way to learn the distribution. The
above discussion has hopefully shown that this claim is false. However,
a related claim, also argued in \Cref{sub:min_entropy} because of its
equivalence in the situations we had then considered, might still be
true: finding the distribution with minimal \textit{coding rate} is a
way to learn the distribution. Supposing for the moment that \(\vx\)
has a density \(p\), the coding rate (i.e., cross-entropy) of the
empirical distribution \(p^{\vX}\) with respect to \(p\) is
\(\CE(p \mmid p^{\vX}) = \infty\). This coding rate is obviously not
minimal in any sense; if we use the same successively better
approximations \(p^{n}\), \(p^{e}\), and \(\hat{q}\) discussed in
\Cref{sub:min_entropy}, it holds
\begin{equation}
    \infty = \CE(p \mmid p^{\vX}) > \CE(p \mmid p^{n}) > \CE(p \mmid
    p^{e}) > \CE(p \mmid \hat{q}) \approx h(p).
\end{equation}
Minimizing the coding rate therefore seems a sensible way to learn the
distribution while avoiding this issue. However, a problem of
tractability now arises: how can we devise algorithms that learn the
distribution by minimizing the coding rate while avoiding the drawbacks
associated with overly minimizing entropy?

One clue is to examine what happens with successfully trained
diffusion models that do not memorize. By writing the trained
denoiser \(\bar{\vx}_{\theta}\) as a perturbation of the memorizing
denoiser \(\bar{\vx}^{\star}\), i.e.,
\begin{equation}
    \bar{\vx}_{\theta}(t, \vx_{t}) = \bar{\vx}^{\star}(t, \vx_{t}) +
    \vdelta_{\theta}(t, \vx_{t})
\end{equation}
we can gain graphical intuition about the action of a well-trained
denoiser, as in \Cref{fig:diffusion_perturbation}. At a very high
level, the trained denoiser computes the memorizing denoiser plus some
perturbation. At a small time \(t\) (perhaps even at the last step of
denoising), the memorizing denoiser points entirely in the direction
of a single training sample (the training sample to be re-generated
by the current iterate). The distribution being learned, or denoised
against, is just the empirical distribution with some
\(\epsilon\)-balls around each point,\footnote{More formally, this
    perturbed distribution has support on (a subset of) a union of
    \(\epsilon\)-balls centered at each training sample. This picture is
supported by past work, e.g., \cite{song2025selective}.} where
\(\epsilon\) is an upper bound on the magnitude of the perturbation,
visualized in \Cref{fig:diffusion_perturbation} by the faint orange
balls.\footnote{In high dimensions, it is hypothesized that a better
    model for these perturbations is not being contained by balls but
    rather by much more anisotropic geometry-adaptive shapes
\cite{farghly2025diffusion}, but the intuition remains similar.} For
increasingly larger well-trained diffusion models, the
perturbation radius \(\epsilon\) shrinks to \(0\), and the model
learns the empirical distribution. On the other hand, \(\epsilon\)
large enough that balls corresponding to different points overlap
or are close may imply that the diffusion model generalizes.

Thus, our successful diffusion model can be viewed as minimizing a
\textit{lossy coding rate}, i.e., coding rate up to a precision
\(\epsilon\) (so as to be able to disregard the \(\epsilon\)-balls
around each point). This notion of the lossy coding rate turns out to
be the right thing to minimize, morally speaking, and will be
\textit{extremely} useful. We will tease out its uses in
\Cref{ch:lossy-compression} and throughout the rest of this book.

\section{Summary and Notes}
\paragraph{Key messages.}
In this Chapter, we have shown that, to pursue a general
low-dimensional distribution, a universal computational approach is
to start with a generic distribution and compress it toward the
distribution through (iterative) noise or entropy reduction. As we
discussed in the preceding Chapter, this is even the case for
arguably the simplest distributions of a single subspace, such as
computing PCA via the power-iteration algorithm. As we will see in
future chapters, the fundamental roles of modern deep networks are
not to fit arbitrary (noisy) data distributions or functions;
instead, they are to realize or approximate those (lossy) compressing
or denoising operators. We will derive precise forms of those
operators and the associated deep architectures in future chapters.

\paragraph{More extensions and references.}
The use of denoising and diffusion for sampling has a rich history.
The first work clearly about a diffusion model is probably
\cite{Sohl-Dickstein2015}, but before this there are many works on
denoising as a computational and statistical problem. The most
relevant of these is probably \cite{hyvarinen05a}, which explicitly
uses the score function to denoise and to perform independent
component analysis. The most popular follow-ups are essentially
co-occurring: \cite{ho2020denoising,song2019}. Since then, thousands
of papers have built on diffusion models; we revisit this topic in
\Cref{ch:consistent-representations}.

Many of these works use a different stochastic process than the
simple linear combination
\eqref{eq:gen_additive_gaussian_noise_model}. In fact, all works
listed above emphasize the need to add \textit{independent} Gaussian
noise at the beginning of each step of the forward process.
Theoretically-minded work uses Brownian motion or stochastic
differential equations to formulate the forward process
\cite{song2020score}. However, since linear combinations of Gaussians
still result in Gaussians, the \textit{marginal distributions} of
such processes still take the form of
\eqref{eq:gen_additive_gaussian_noise_model}. Most of our discussion
requires only that the marginal distributions are what they are, so
our overly simplistic model is actually enough for almost everything.
The only time marginal distributions are not enough is when we derive
an expression for \(\Ex[\vx_{s} \mid \vx_{t}]\) in terms of
\(\Ex[\vx \mid \vx_{t}]\). Different (noising) processes give
different such expressions, which can be used for sampling (and of
    course there are other ways to derive efficient samplers, such as the
ever-popular DDPM sampler). The process in
\eqref{eq:gen_additive_gaussian_noise_model} is a bona fide
stochastic process whose ``natural'' denoising iteration takes the
form of the popular DDIM algorithm
\cite{song2020denoising}. (Even this equivalence is not trivial; we
cite \cite{de2025distributional} as justification.)

On top of the theoretical work \citep{li2024d} covered in
\Cref{sub:sampling_denoising}, and the lineage of work that it builds
on, which studies the \textit{sampling} efficiency of diffusion
models when the data has low-dimensional structure, there is a large
body of work that studies the \textit{training} efficiency of
diffusion models when the data has low-dimensional structure.
Specifically, \citet{chen2023score} and \citet{wang2024diffusion}
characterized the approximation and estimation error of denoisers
when the data belongs to a mixture of low-rank Gaussians, showing
that the number of training samples required to accurately learn the
distribution scales with the intrinsic dimension of the
data rather than the ambient dimension. There is considerable
\textit{methodological} work that attempts to use the
low-dimensional structure of the data to perform various tasks with
diffusion models. We highlight three: image editing
\citep{chen2024exploring}, watermarking \citep{li2024shallow}, and
unlearning \citep{chen2025dual}, though this list is inexhaustive.

\section{Exercises and Extensions}

\begin{exercise}
    Show that \eqref{eq:optimal_denoiser} is the optimal solution
    of Problem \eqref{eq:denoising_loss}.
\end{exercise}

\begin{exercise}\label{exercise:conditional_gaussian}
    Consider random vectors $\vx \in \bR^D$ and $\vy \in \bR^d$,
    such that the
    pair $(\vx, \vy) \in \bR^{D + d}$ is jointly Gaussian. This means that
    \begin{equation*}
        \begin{bmatrix}
            \vx \\
            \vy
        \end{bmatrix}
        \sim
        \cN \left(
            \begin{bmatrix}
                \vmu_{\vx} \\
                \vmu_{\vy}
            \end{bmatrix}
            ,
            \begin{bmatrix}
                \vSigma_{\vx} & \vSigma_{\vx\vy} \\
                \vSigma_{\vx\vy}^\top & \vSigma_{\vy}
            \end{bmatrix}
        \right),
    \end{equation*}
    where the mean and covariance parameters are given by
    \begin{equation*}
        \vmu_{\vx} = \bE[\vx],\quad \vmu_{\vy} = \bE[\vy],\quad
        \begin{bmatrix}
            \vSigma_{\vx} & \vSigma_{\vx\vy} \\
            \vSigma_{\vx\vy}^\top & \vSigma_{\vy}
        \end{bmatrix}
        =
        \bE\left[
            \begin{bmatrix}
                \vx - \bE[\vx] \\
                \vy - \bE[\vy]
            \end{bmatrix}
            \begin{bmatrix}
                \vx - \bE[\vx] \\
                \vy - \bE[\vy]
            \end{bmatrix}^\top
        \right]
    \end{equation*}
    Assume that $\vSigma_{\vy}$ is positive definite (hence
    invertible); then
    positive semidefiniteness of the covariance matrix is equivalent
    to the Schur
    complement condition $\vSigma_{\vx} - \vSigma_{\vx\vy}
    \vSigma_{\vy}^{-1}
    \vSigma_{\vx\vy}^\top \succeq \Zero$.

    In this exercise, we will prove that the conditional distribution
    $p_{\vx \mid
    \vy}$ is Gaussian: namely,
    \begin{equation}\label{eq:gaussian-conditional-eqn}
        p_{\vx \mid \vy} \sim \cN\left(
            \vmu_{\vx} + \vSigma_{\vx\vy} \vSigma_{\vy}^{-1} (\vy
            - \vmu_{\vy}),
            \vSigma_{\vx} - \vSigma_{\vx\vy} \vSigma_{\vy}^{-1}
            \vSigma_{\vx\vy}^{\top}
        \right).
    \end{equation}
    A direct path to prove this result manipulates the defining ratio of
    densities $p_{\vx, \vy} / p_{\vy}$. We sketch an algebraically concise
    argument of this form below.

    \begin{enumerate}
        \item Verify the following matrix identity for the covariance:
            \begin{equation}\label{eq:gaussian-conditional-covariance-block}
                \begin{bmatrix}
                    \vSigma_{\vx} & \vSigma_{\vx\vy} \\
                    \vSigma_{\vx\vy}^\top & \vSigma_{\vy}
                \end{bmatrix}
                =
                \begin{bmatrix}
                    \vI_D & \vSigma_{\vx\vy}\vSigma_{\vy}^{-1} \\
                    \Zero & \vI_d
                \end{bmatrix}
                \begin{bmatrix}
                    \vSigma_{\vx} - \vSigma_{\vx\vy} \vSigma_{\vy}^{-1}
                    \vSigma_{\vx\vy}^{\top} & \Zero \\
                    \Zero & \vSigma_{\vy}
                \end{bmatrix}
                \begin{bmatrix}
                    \vI_D & \Zero\\
                    \vSigma_{\vy}^{-1}\vSigma_{\vx\vy}^\top & \vI_d
                \end{bmatrix}.
            \end{equation}
            One arrives at this identity by performing two rounds of
            (block) Gaussian
            elimination on the covariance matrix.
        \item Based on the previous identity, show that
            \begin{equation}\label{eq:gaussian-conditional-covariance-inverse-block}
                \begin{bmatrix}
                    \vSigma_{\vx} & \vSigma_{\vx\vy} \\
                    \vSigma_{\vx\vy}^\top & \vSigma_{\vy}
                \end{bmatrix}^{-1}
                =
                \vP
                \begin{bmatrix}
                    \left(\vSigma_{\vx} - \vSigma_{\vx\vy}
                        \vSigma_{\vy}^{-1}
                    \vSigma_{\vx\vy}^{\top}\right)^{-1} & \Zero \\
                    \Zero & \vSigma_{\vy}^{-1}
                \end{bmatrix}
                \vP^\top,
            \end{equation}
            where
            \begin{equation}
                \vP =
                \begin{bmatrix}
                    \vI_D & \Zero\\
                    -\vSigma_{\vy}^{-1}\vSigma_{\vx\vy}^\top & \vI_d
                \end{bmatrix},
            \end{equation}
            and whenever the relevant inverses are
            defined.\footnote{In cases where the
                Schur complement term is not invertible, the same
                result holds with its
                inverse replaced by the Moore-Penrose pseudoinverse.
                In particular, the
                conditional distribution
                \eqref{eq:gaussian-conditional-eqn} becomes
            a degenerate Gaussian distribution.}
            Conclude that
            \begin{align}
                &
                \begin{bmatrix}
                    \vx-\vmu_{\vx} \\
                    \vy-\vmu_{\vy}
                \end{bmatrix}^{\top}
                \begin{bmatrix}
                    \vSigma_{\vx} & \vSigma_{\vx\vy} \\
                    \vSigma_{\vx\vy}^\top & \vSigma_{\vy}
                \end{bmatrix}^{-1}
                \begin{bmatrix}
                    \vx-\vmu_{\vx} \\
                    \vy-\vmu_{\vy}
                \end{bmatrix}
                \\
                &\qquad=
                \begin{bmatrix}
                    \vx - \vm \\
                    \vy - \vmu_{\vy}
                \end{bmatrix}^\top
                \begin{bmatrix}
                    \left(\vSigma_{\vx} - \vSigma_{\vx\vy}
                        \vSigma_{\vy}^{-1}
                    \vSigma_{\vx\vy}^{\top}\right)^{-1} & \Zero \\
                    \Zero & \vSigma_{\vy}^{-1}
                \end{bmatrix}
                \begin{bmatrix}
                    \vx - \vm \\
                    \vy - \vmu_{\vy}
                \end{bmatrix},
            \end{align}
            where $\vm = \vmu_{\vx} +
            \vSigma_{\vx\vy}\vSigma_{\vy}^{-1}(\vy
            - \vmu_{\vy})$.
            %(\textit{Hint: To economize algebraic manipulations, note
            % that the first
            %and last matrices on the RHS of
            %\Cref{eq:gaussian-conditional-covariance-block} are
            % transposes of one
            %another.})
        \item By dividing $p_{\vx, \vy} / p_{\vy}$, prove
            \Cref{eq:gaussian-conditional-eqn}. (\textit{Hint:
                    Using the previous
                    identities, only minimal algebra should
                    be necessary.
                    For the normalizing
                    constant, use
                    \Cref{eq:gaussian-conditional-covariance-inverse-block}
                    to
            factor the determinant similarly.})

    \end{enumerate}

\end{exercise}

\begin{exercise}\label{exercise:sherman_morrison_woodbury_identity}
    Show the Sherman--Morrison--Woodbury identity, i.e., for matrices
    \(\vA\), \(\vC\), \(\vU\), \(\vV\) such that \(\vA\), \(\vC\), and
    \(\vA + \vU\vC\vV\) are invertible,
    \begin{equation}
        (\vA + \vU\vC\vV)^{-1} = \vA^{-1} - \vA^{-1}\vU(\vC^{-1} +
        \vV\vA^{-1}\vU)^{-1}\vV\vA^{-1}.
    \end{equation}
\end{exercise}

\begin{exercise}\label{exercise:generalizing_results_to_different_noise_models}
    Rederive the following, assuming \(\vx_{t}\) follows the generalized
    noise model \eqref{eq:gen_additive_gaussian_noise_model}.
    \begin{itemize}
        \item Tweedie's formula: \eqref{eq:gen_tweedie}.
        \item The DDIM iteration: \eqref{eq:gen_denoising_iteration}.
        \item The Bayes-optimal denoiser for a Gaussian mixture model:
            \eqref{eq:gen_gmm_bayes_optimal_denoiser}.
    \end{itemize}
\end{exercise}

\begin{exercise}\label{exercise:implement_denoising_processes}
    \phantom{}
    \begin{enumerate}
        \item Implement the formulae derived in
            \Cref{exercise:generalizing_results_to_different_noise_models},
            building a sampler for Gaussian mixtures.
        \item Reproduce \Cref{fig:ve_forward_denoising} and
            \Cref{fig:vp_gmm_denoising}.
        \item We now introduce \textit{Flow
            Matching (FM)} (see \Cref{rem:flow_matching}), as follows:
            \begin{equation}
                \alpha_{t} = 1 - t, \qquad \sigma_{t} = t.
            \end{equation}
            Implement this process using the same framework, and test it
            for sampling in high dimensions. Which process gives better
            or more stable results?
    \end{enumerate}
\end{exercise}

\begin{exercise}\label{exercise:velocity_prediction}
    Consider the somewhat familiar interpolation
    \begin{equation}
        \vx_{t} = \alpha_{t}\vx + \sigma_{t}\vy, \qquad \forall t
        \in [0, T].
    \end{equation}
    where \(\vy\) is a (not necessarily Gaussian) random
    variable. We define the \textit{velocity} as
    \begin{equation}
        \vv_{t} = \odv*{\vx_{t}}{t} = \alpha_{t}^{\prime}\vx +
        \sigma_{t}^{\prime}\vy, \qquad \forall t \in [0, T].
    \end{equation}
    Compute an algebraic relationship writing \(\Ex[\vv_{t} \mid
    \vx_{t}]\) in terms of \(\Ex[\vx \mid \vx_{t}]\) and vice
    versa, showing that velocity prediction is algebraically
    equivalent to denoising and thus also to noise prediction
    (see \Cref{rem:flow_matching}).
\end{exercise}

\end{document}
